\documentclass[a4paper,10pt]{article}
\usepackage[multidot]{grffile}
\usepackage{fullpage}
\usepackage{graphicx}
\usepackage{txfonts}
\usepackage{natbib}
\usepackage[colorlinks,allcolors=blue,bookmarks=false,hypertexnames=true]{hyperref} 
\hyphenation{get-sources get-filaments get-images get-old non-uniform}
\catcode`\_=11\relax
\newcommand\email[1]{\_email #1\q_nil}
\def\_email#1@#2\q_nil{%
  \href{mailto:#1@#2}{{\emailfont\detokenize{#1}\emailampersat\detokenize{#2}}}%
}
\newcommand\emailfont{\sffamily}
\newcommand\emailampersat{{\color{red}\small@}}
\catcode`\_=8\relax   
\newcommand*{\arcsec}{^{\prime\prime}\mkern-1.2mu}
\usepackage{verbatimbox}
\usepackage{latexsym}
\usepackage[english]{isodate}
\pdfoutput=1

%%\def\mydate{\leavevmode\hbox{\the\year/\twodigits\month/\twodigits\day}}
%%\def\twodigits#1{\ifnum#1<10 0\fi\the#1}

\title
{
\textbf{User's Guide\\Extracting Sources and Filaments with \textsl{getsf}}
}

\author
{
Alexander Men'shchikov\vspace{1.2mm}\\
\normalsize
AIM, CEA, CNRS, Universit{\'e} Paris-Saclay, Universit{\'e} Paris Diderot, Sorbonne Paris Cit{\'e}\\
\normalsize
F-91191 Gif-sur-Yvette, France
}

\date{} % <-- This is here to remove the automatic date insertion by \maketitle
\begin{document}
\null  % Empty line
\nointerlineskip  % No skip for prev line
\vfill
\let\snewpage \newpage
\let\newpage \relax
\vspace{-30mm}
\maketitle
\vspace{-3mm}
\centerline{\date{\today}}
\vspace{5mm}
\centerline{{The \textsl{getsf} method was published by \cite{Men'shchikov2021getsf}}}
\vspace{3mm}
\centerline{\normalsize{\url{http://irfu.cea.fr/Pisp/alexander.menshchikov/}}}
\centerline{\normalsize{\url{https://ascl.net/2012.001}}}
\thispagestyle{empty}  % No page number
\let \newpage \snewpage
\vfill 
\break % page break
%\maketitle

\newpage
\tableofcontents
\newpage

%||||||||||||||||||||||||||||||||||||||||||||||||||||||||||||||||||||||||||||||||||||||||||||||||||||||||||||||||||||||||||||||||||
%||||||||||||||||||||||||||||||||||||||||||||||||||||||||||||||||||||||||||||||||||||||||||||||||||||||||||||||||||||||||||||||||||

\section{Introduction}
\label{intro}

The term \emph{source} denotes almost round intensity emission peaks in the observed images, significantly stronger than the local
surrounding fluctuations, indicating the presence of one or several physical \emph{objects} in space that produced the emission. An
implicit assumption that an unresolved far-infrared source on a complex fluctuating background contains emission of just one single
object is invalid in general: the unresolved peak of emission may well be a blend of several unrelated components, produced by
different physical entities. The term \emph{filament} refers to significantly elongated structures of various shapes. The term
\emph{background} refers to the approximately isotropic astrophysical backgrounds, such as molecular clouds, emission nebulae, or
H\textsc{ii} regions. The term \emph{noise} refers to the instrumental noise with possible contributions from any other signals that
are not astrophysical in nature, i.e., unrelated to the emission of the areas in space being observed.

Like other extraction algorithms, \textsl{getsf} does not know anything about the physical nature of the objects that produce the
detected sources of emission. It extracts \emph{sources}, separating them from background, filaments, and noise. It is a subject of
the post-extraction investigation to identify and sub-select from the catalogs those sources of infrared emission that are produced
by the physical objects of a specific nature.
     
\begin{quote} %--------------------------------------------------------------------------------------------------------------------
\textbf{NOTE}. The \textsl{getsf} method has no free parameters to vary. It has been carefully tested and fine-tuned to produce the
best results with the default values of the configuration parameters and the common user should not change them. There is only one
parameter (per image) that a user must specify before any extraction: the maximum size (FWHM) of the sources or filaments of
interest \citep[see Sect.~3.1.3 in][]{Men'shchikov2021getsf}. This can easily be done by opening an image in \textsl{ds9} (or 
another image viewer) and placing a region fully encircling the \emph{entire} emission peak of the largest structure; radius of the 
region (in arcsec) would be the maximum size.
\end{quote} %----------------------------------------------------------------------------------------------------------------------

%||||||||||||||||||||||||||||||||||||||||||||||||||||||||||||||||||||||||||||||||||||||||||||||||||||||||||||||||||||||||||||||||||
%||||||||||||||||||||||||||||||||||||||||||||||||||||||||||||||||||||||||||||||||||||||||||||||||||||||||||||||||||||||||||||||||||

\section{The \textsl{getsf} software suite}
\label{software}

The tool has been developed as a \textsl{bash} script \textsl{getsf} that employs several utilities (in FORTRAN) to perform all
numerical work. The Linux or macOS operating systems have been used to install and run the code.

%||||||||||||||||||||||||||||||||||||||||||||||||||||||||||||||||||||||||||||||||||||||||||||||||||||||||||||||||||||||||||||||||||

\subsection{Prerequisites}
\label{prerequisites}

The \textsl{getsf} software can be installed on any modern laptop, desktop, or server. For large images, it requires many gigabytes
of memory and space, hence the more the better. It is necessary to verify that a FORTRAN compiler, either
\textsl{ifort}\footnote{\url{https://software.intel.com/content/www/us/en/develop/tools/oneapi/components/fortran-compiler.html}}
or \textsl{gfortran}\footnote{\url{https://github.com/fxcoudert/gfortran-for-macOS/releases}}, is available in the operating system:
%----------------------------------------------------------------------------------------------------------------------------------
\begin{verbnobox}[\mbox{}\small] 
> ifort --version
> gfortran --version
\end{verbnobox} 
%----------------------------------------------------------------------------------------------------------------------------------
If neither is found, one of them must be installed. For reading and writing FITS images, \textsl{getsf} uses the \textsl{cfitsio}
library \citep{Pence1999}, hence \verb!libcfitsio.a! must exist in your operating system. For resampling and reprojecting images,
\textsl{getsf} calls \textsl{swarp} \citep{Bertin_etal2002}, hence it is necessary to verify its availability:
%----------------------------------------------------------------------------------------------------------------------------------
\begin{verbnobox}[\mbox{}\small] 
> swarp --version
\end{verbnobox} 
%----------------------------------------------------------------------------------------------------------------------------------
If \textsl{swarp} is not found, it needs to be installed (\url{http://www.astromatic.net/software/swarp/}). For determining the
world coordinates of sources, \textsl{getsf} employs \textsl{xy2sky} from \textsl{wcstools} \citep{Mink2002}:
%----------------------------------------------------------------------------------------------------------------------------------
\begin{verbnobox}[\mbox{}\small] 
> xy2sky -version
\end{verbnobox} 
%----------------------------------------------------------------------------------------------------------------------------------
If \textsl{xy2sky} is not found, it needs to be installed (\url{http://tdc-www.harvard.edu/wcstools/}) to ensure that the source
catalogs have world coordinates. Without this (optional) utility, the source catalogs would have coordinates in pixels only.
Finally, to make the \textsl{getsf} screen output colored, it is recommended to have the \textsl{highlight} utility:
%----------------------------------------------------------------------------------------------------------------------------------
\begin{verbnobox}[\mbox{}\small] 
> highlight --version
\end{verbnobox} 
%----------------------------------------------------------------------------------------------------------------------------------
This utility is optional (\url{http://www.andre-simon.de/}), it just improves visibility of the screen output.

%||||||||||||||||||||||||||||||||||||||||||||||||||||||||||||||||||||||||||||||||||||||||||||||||||||||||||||||||||||||||||||||||||

\subsection{Installation of \textsl{getsf}}
\label{install}

The software is compressed in a \textsl{zip} archive \verb!getsf.vYYMMDD.zip!, whose name contains the latest \textsl{getsf}
version (e.g., \verb!getsf.v241122.zip!). It is recommended that the downloaded archive is moved to a dedicated directory (for
example, \verb!/Users/yourname/+GETSF!), where it must be unpacked:
%----------------------------------------------------------------------------------------------------------------------------------
\begin{verbnobox}[\mbox{}\small] 
> unzip getsf.v241122.zip
\end{verbnobox} 
%----------------------------------------------------------------------------------------------------------------------------------
The unpacked code is created in \verb!/Users/yourname/+GETSF/v241122!. It is necessary to add the environment variable
\verb!GETSF_HOME! to the $\sim$\verb!/.bashrc! file:
%----------------------------------------------------------------------------------------------------------------------------------
\begin{verbnobox}[\mbox{}\small] 
export GETSF_HOME=/Users/yourname/+GETSF
\end{verbnobox} 
%----------------------------------------------------------------------------------------------------------------------------------
It is also necessary to create the binary directory, where the executable code will be found, and add the environment variable
\verb!GETSF_BIN! to the $\sim$\verb!/.bashrc! file:
%----------------------------------------------------------------------------------------------------------------------------------
\begin{verbnobox}[\mbox{}\small] 
> mkdir -p /Users/yourname/bin/+GETSF
export GETSF_BIN=/Users/yourname/bin/+GETSF
\end{verbnobox} 
%----------------------------------------------------------------------------------------------------------------------------------
The installation script \textsl{installg} will put the compiled binaries and scripts in \verb!GETSF_BIN!. Additionally, this
current version will also be installed in a dedicated directory \verb!/Users/yourname/bin/+GETSF/v241122!, which allows running
both the default and specific versions of the code, when necessary. The \verb!PATH! environment variable in the
$\sim$\verb!/.bashrc! file must be modified to include the binary directory, for example:
%----------------------------------------------------------------------------------------------------------------------------------
\begin{verbnobox}[\mbox{}\small] 
export PATH=$GETSF_BIN:$PATH
\end{verbnobox} 
%----------------------------------------------------------------------------------------------------------------------------------
The \textsl{cfitsio} library \verb!libcfitsio.a! must be copied from its default location in your operating system to the
\textsl{getsf} home directory, for example:
%----------------------------------------------------------------------------------------------------------------------------------
\begin{verbnobox}[\mbox{}\small] 
> cp /path/to/cfitsio/library/in/your/system/libcfitsio.a "$GETSF_HOME"
\end{verbnobox} 
%----------------------------------------------------------------------------------------------------------------------------------
The \textsl{getsf} installation must be started from its directory
\verb!$GETSF_HOME/v241122!:
%----------------------------------------------------------------------------------------------------------------------------------
\begin{verbnobox}[\mbox{}\small] 
> ./installg
\end{verbnobox} 
%----------------------------------------------------------------------------------------------------------------------------------
The installation script compiles all utilities and copies their binaries and other scripts in the two directories:\\
%----------------------------------------------------------------------------------------------------------------------------------
\vspace{-2mm}\noindent\\
{\small\verb!> $GETSF_BIN!} $\,\,\leadsto$ binary directory of the default version\\
{\small\verb!> $GETSF_BIN/v241122!\vspace{-2mm}} $\,\,\leadsto$ binary directory of the specific version\\\\
%----------------------------------------------------------------------------------------------------------------------------------
To make sure that the installed \textsl{getsf} is indeed available, its utilities may be run, for example:
%----------------------------------------------------------------------------------------------------------------------------------
\begin{verbnobox}[\mbox{}\small] 
> modfits
> operate
\end{verbnobox} 
%----------------------------------------------------------------------------------------------------------------------------------
The default version of \textsl{getsf} can be started by typing its name:
%----------------------------------------------------------------------------------------------------------------------------------
\begin{verbnobox}[\mbox{}\small] 
> getsf
\end{verbnobox} 
%----------------------------------------------------------------------------------------------------------------------------------
The specific version of \textsl{getsf} can be started by typing its name and the version number:
%----------------------------------------------------------------------------------------------------------------------------------
\begin{verbnobox}[\mbox{}\small] 
> getsf 241122
\end{verbnobox} 
%----------------------------------------------------------------------------------------------------------------------------------
Upon each \textsl{getsf} run, several log files are created and/or appended:\\
%----------------------------------------------------------------------------------------------------------------------------------
\vspace{-2mm}\noindent\\
{\small\verb!> +log.AGREE!} $\,\,\leadsto$ user agreement and suggested citations\\
{\small\verb!> +log.HISTORY!} $\,\,\leadsto$ history of recent changes in \textsl{getsf} and its utilities\\
{\small\verb!> +log.VERSION!} $\,\,\leadsto$ brief information on the version of the code executed\\
{\small\verb!> +log.USAGE!} $\,\,\leadsto$ brief \textsl{getsf} usage information\\
{\small\verb!> +log.GETSF!} $\,\,\leadsto$ main log file with information on the extraction process\\
{\small\verb!> +log.GETSF.iter!} $\,\,\leadsto$ iterations log for the single-scale image processing\\
{\small\verb!> +log.FOOTS.iter!\vspace{-2mm}} $\,\,\leadsto$ iterations log for the source footprints in measurements\\\\
%----------------------------------------------------------------------------------------------------------------------------------

\begin{quote} %--------------------------------------------------------------------------------------------------------------------
\textbf{NOTE}. The software may become corrected and improved in the future and only the most recent version of \textsl{getsf} is
supported. If apparently abnormal results are found or the code aborts for any reason, the user is asked to install the latest
version of the code from its website and restart the failed action to see if the problem persists. If it does, please refer to
\url{http://irfu.cea.fr/Pisp/alexander.menshchikov/#support} and request assistance from the author.
\end{quote} %----------------------------------------------------------------------------------------------------------------------

%||||||||||||||||||||||||||||||||||||||||||||||||||||||||||||||||||||||||||||||||||||||||||||||||||||||||||||||||||||||||||||||||||

\subsection{The utilities and scripts}
\label{utils}

The following list of the \textsl{getsf} utilities and scripts explains their purpose and functions. They are quite useful for
command-line image manipulations, even if there is no need in the source or filament extraction. The utilities are sorted in a
decreasing sequence of their general usability outside \textsl{getsf}.
\begin{tabbing} %------------------------------------------------------------------------------------------------------------------
Utility \,\,\,\,\,\,\,\, \= Purpose \kill
\textsl{modfits}   \> $\leadsto$ modify an image or its header in various ways: math transformations; profiling an image along a\\
                   \> \;\;\:\: line; image segmentation; filament skeletonization; removal of connected clusters of pixels; addition\\
                   \> \;\;\:\: or removal of border areas; correction of saturated or bad pixel areas; conversion of intensity units;\\
                   \> \;\;\:\: changes of header keywords; \emph{etc}.\\
\textsl{operate}   \> $\leadsto$ operate on two input images: addition, subtraction, multiplication, division; relative differencing;\\
                   \> \;\;\:\: minimization or maximization; extension or expansion of masks; copying of an image header;\\
                   \> \;\;\:\: computation of surface densities, temperatures, or intensities; \emph{etc}.\\
\textsl{imgstat}   \> $\leadsto$ display and/or save image statistical quantities; produce mode-, mean-, or median-filtered images;\\
                   \> \;\;\:\: compute images of standard deviations, skewness, kurtosis; \emph{etc}.\\
\textsl{fftconv}   \> $\leadsto$ fast Fourier transform or convolve image with few predefined kernels or an external kernel image\\
\textsl{fitfluxes} \> $\leadsto$ fit spectral shapes of source fluxes or image pixel intensities to derive masses or surface densities\\
\textsl{readhead}  \> $\leadsto$ display an image header or save selected keywords\\
\textsl{convolve}  \> $\leadsto$ convolve an image to a desired lower resolution\\
\textsl{resample}  \> $\leadsto$ resample and re-project an image with optional rotation\\
\textsl{fmeasure}  \> $\leadsto$ measure and catalog properties of detected filaments\\
\textsl{fixskels}  \> $\leadsto$ recreate scale-dependent skeletons as the isolated ones to produce global skeletons\\
\textsl{hires}     \> $\leadsto$ derive high-resolution surface densities and temperatures\\
\textsl{prepobs}   \> $\leadsto$ convert observed images to the same pixel grid\\
\textsl{installg}  \> $\leadsto$ install \textsl{getsf} in an operating system (macOS, Linux)\\
\textsl{iospeed}   \> $\leadsto$ test I/O speed of a hard drive for a specific image\\
\textsl{cleanbg}   \> $\leadsto$ interpolate background under footprints of sources\\
\textsl{ellipses}  \> $\leadsto$ overlay an image with ellipses of extracted sources\\
\textsl{sfinder}   \> $\leadsto$ detect sources in combined single-scale images\\
\textsl{smeasure}  \> $\leadsto$ measure and catalog properties of detected sources\\
\textsl{finalcat}  \> $\leadsto$ produce final catalogs of detected sources\\
\textsl{expanda}   \> $\leadsto$ expand masked areas of an image to its edges\\
\textsl{extractx}  \> $\leadsto$ extract all FITS extensions in separate images\\
\textsl{splitcube} \> $\leadsto$ split a FITS data cube in separate images\\
\textsl{makecube}  \> $\leadsto$ make a FITS data cube from separate images
\end{tabbing} %--------------------------------------------------------------------------------------------------------------------
\begin{quote} %--------------------------------------------------------------------------------------------------------------------
\textbf{NOTE}. With \textsl{getsf} utilities, it is not necessary to supply the \verb!.fits! extension of the image name on the
command line. When a utility is run without any parameter, its usage information is displayed.
\end{quote} %----------------------------------------------------------------------------------------------------------------------

%||||||||||||||||||||||||||||||||||||||||||||||||||||||||||||||||||||||||||||||||||||||||||||||||||||||||||||||||||||||||||||||||||

\subsection{Notable improvements or new features}
\label{improvements}

\begin{itemize} 

\item Revision \verb!260706! (improvement). This version of \emph{getsf} substantially improves interpolated backgrounds of 
sources. Instead of a simple interpolation along 4 directions used previously, \emph{cleanbg} now interpolates along arbitrary 
number of directions, defaulting to 16. This improves source backgrounds, eliminating linear artifacts that sometimes appeared 
in complex surroundings with the simpler algorithm. This version also corrects \emph{tools} functions \emph{fftconvolve}, 
\emph{rlft3}, and \emph{fourn} to eliminate very low-level stripes in convolved images.

\item Revision \verb!250609! (new feature). This version of \textsl{getsf} enables filament measurements on all spatial scales.
Previously, measurements were done only on the scale of the maximum width of the filaments of interest, defined in the
\verb!+getsf.cfg! (Sect.~\ref{params}), because background of filaments was derived for that scale only, hence the background
subtraction was not accurate for much narrower filaments. To derive additional backgrounds on smaller scales and obtain accurate
measurements, it is necessary to make several short runs of \textsl{getsf} for a few processing steps, setting smaller values of
the maximum width of the filaments. For more details, users are referred to Sect.~\ref{measureall}.

\item Revision \verb!250603! (new feature). Starting with this version, \textsl{getsf} includes a new utility \textsl{fixskels}
that recreates scale-dependent skeletons as the isolated ones -- not overlapping with skeletons on any of the larger scales. The
global skeletons are now correctly obtained by a simple summation of the isolated scale-dependent skeletons over all scales. The
new global skeletons provide most accurate visualization of detected filaments on all spatial scales. They can also be used for
filament measurements independently (outside \textsl{getsf}), if a user finds it useful. In this version, structure of the
subdirectories for the resulting skeletons and filaments was slightly rearranged, as well as some output files were renamed to
better reflect their contents.

\item Revision \verb!241206! (improvement). This version improves \textsl{fmeasure} in production of filament footprints and
corrects a problem in calculation of the filament masses and corresponding linear densities (mass/length) over the one-sided 
footprints. Masses and linear densities for some blended filaments may change, as a result. Images of the one-sided footprints are 
no more created.

\item Revision \verb!241122! (improvement). Starting with this version, \textsl{getsf} produces filament measurements for only
\emph{three} spatial scales centered at the maximum width of the filaments of interest, set in the \textsl{getsf} configuration
file \verb!+getsf.cfg! (Sect.~\ref{params}). This is aimed at avoiding confusion, because the meaning of the maximum width of
filaments is somewhat different from that of the maximum size of sources, in the context of filament measurements. The background
of filaments is determined for the filaments detected on spatial scales around the adopted maximum width, therefore subtraction of
the background is valid only for those filaments. Measurements of the filaments detected on significantly different scales are
likely to be inaccurate, because of the inappropriate background subtraction. To properly measure also the narrower (or wider)
filaments detected on a smaller (or larger) spatial scale, \textsl{getsf} must be additionally run using the corresponding smaller
(or larger) value of the maximum width parameter in \verb!+getsf.cfg! (Sect.~\ref{params}).

\item Revision \verb!220627! (new feature). Starting with this version, \textsl{getsf} includes \textsl{makecube}, a new utility
that allows creation of 3D data cubes from separate 2D images to facilitate analysis of extraction results, primarily when running
\textsl{getsf} on slices extracted from line data cubes.

\item Revision \verb!220610! (new feature). Starting with this version, \textsl{getsf} creates \textsl{ds9} region files for both
sources (half-maximum and footprint ellipses) and filaments (global and scale-dependent skeletons) to facilitate their visualization
and analyses. The region files \verb!*.reg! are found in relevant subdirectories (Sects.~\ref{srccats}, \ref{srcvisuals}, 
\ref{filprofiles}, and \ref{filvisuals}).

%\item Revision \verb!210915! (new feature). Starting with this version, \textsl{getsf} provides also the measurements of filaments
%for the global skeletons the same way the measurements are done for the set of scale-dependent skeletons. The images and catalogs
%\verb!*.ske.global.*! are placed in the usual subdirectories (Sect.~\ref{filprofiles}).

\item Revision \verb!210907! (new feature). Starting with this version, \textsl{getsf} adds the \emph{global}, scale-independent
skeletons to its output. Several users experienced difficulties analyzing the complex nature of observed filaments with the set of
scale-dependent skeletons created by \textsl{getsf}. In response, the updated \textsl{getsf} produces also images of the global
skeletons tracing the filamentary structures on all scales, independently of the filament width. To emphasize: the global skeletons
are provided for \emph{qualitative visualization and illustration purposes} -- they are not as reliable as the scale-dependent
skeletons, so they should not be used for measurements of the filaments. The global skeletons are approximately created from the
scale-dependent skeletons within a range of scales of interest, up to the maximum width of the filaments specified in the
\textsl{getsf} configuration file (Sect.~\ref{params}). The additional images of the global skeletons \verb!*ske.global.fits!, are
located in the subdirectories \verb!+skeletons! and \verb!+filaments/+visuals! (Sects.~\ref{skeletons} and \ref{filvisuals}).

\end{itemize}

%||||||||||||||||||||||||||||||||||||||||||||||||||||||||||||||||||||||||||||||||||||||||||||||||||||||||||||||||||||||||||||||||||
%||||||||||||||||||||||||||||||||||||||||||||||||||||||||||||||||||||||||||||||||||||||||||||||||||||||||||||||||||||||||||||||||||

\section{Preparing images for extraction}
\label{preparation}

The first step is to create the extraction directories (anywhere in the file system, may be named freely). In practice, it makes
sense to use current date in the extraction directory name (in the \verb!YYMMDD! format) and the name of the observed region, for 
example:\\
%----------------------------------------------------------------------------------------------------------------------------------
\vspace{-2mm}\noindent\\
{\small\verb!mkdir -p images!} $\,\,\leadsto$ preparation directory (name may be arbitrary)\\
{\small\verb!mkdir -p 241122/runs!\vspace{-2mm}} $\,\,\leadsto$ extraction directory (names may be arbitrary)\\\\
%----------------------------------------------------------------------------------------------------------------------------------
This document uses the names from the above example. It is assumed that the observed data set is found in the directory
\verb!images!, which is made the current directory in a terminal window. If the data set consists of more than a single image, the
images usually have different pixel sizes and coverage areas; for use with \textsl{getsf}, all images need to be reprojected to the
same grid of pixels.

%||||||||||||||||||||||||||||||||||||||||||||||||||||||||||||||||||||||||||||||||||||||||||||||||||||||||||||||||||||||||||||||||||

\subsection{Handling multiple extensions and data cubes}
\label{hipecube}

Observed images may come in the FITS files with multiple extensions (e.g., the \emph{Herschel} HIPE output) or in the 3D data
cubes. To use the images with \textsl{getsf}, it is necessary to extract them from the files as plain 2D FITS images. This can
easily be done with the \textsl{getsf} utilities \textsl{extractx} and \textsl{splitcube}:
%----------------------------------------------------------------------------------------------------------------------------------
\begin{verbnobox}[\mbox{}\small]
> extractx hipe.file.fits
> splitcube data.cube.fits 3
\end{verbnobox}
%----------------------------------------------------------------------------------------------------------------------------------
An optional parameter \verb!1|2|3! of \textsl{splitcube} defines the data axis, along which the cube must be split into the normal
2D images, with the default value of \verb!3!. The resulting output images are named as
%----------------------------------------------------------------------------------------------------------------------------------
\begin{verbnobox}[\mbox{}\small]
hipe.file.01.fits, hipe.file.02.fits, hipe.file.03.fits, ...
data.cube.0001.fits, data.cube.0002.fits, data.cube.0003.fits, ...
\end{verbnobox}
%----------------------------------------------------------------------------------------------------------------------------------
In the HIPE files, observed images are usually found in the second FITS extension. The observed position-velocity data cubes
usually contain too many slices, from the practical point of view. To use them with \textsl{getsf} (that handles up to 99 images),
the slices can be averaged over carefully chosen intervals of velocities to reduce the number of images well below 100. For their
first test extractions, the \textsl{getsf} users are advised to have two or three images.

%||||||||||||||||||||||||||||||||||||||||||||||||||||||||||||||||||||||||||||||||||||||||||||||||||||||||||||||||||||||||||||||||||

\subsection{Modifying FITS headers}
\label{headers}

It is not unusual that the headers of FITS images are either incorrect or incomplete. The headers can be displayed and modified
using the \textsl{readhead} and \textsl{modfits} utilities, respectively. The standard minimal subset of the FITS header keywords,
supported by \textsl{getsf} and its utilities, is shown by the following example:
%----------------------------------------------------------------------------------------------------------------------------------
\begin{verbnobox}[\mbox{}\small]
> readhead universe_250

 READHEAD • Display FITS Headers • Nov 05 2020 17:25
 Alexander Men’shchikov, DAp IRFU CEA Saclay, France.
 Using CFITSIO library version 3.260 by William D Pence.

   Reading header of 'universe_250.fits'...

 SIMPLE  =                    T / file does conform to FITS standard
 BITPIX  =                  -32 / number of bits per data pixel
 NAXIS   =                    2 / number of data axes
 NAXIS1  =                 2600 / length of data axis 1
 NAXIS2  =                 3500 / length of data axis 2
 COMMENT   FITS (Flexible Image Transport System) format is defined in 'Astronomy
 COMMENT   and Astrophysics', volume 376, page 359; bibcode: 2001A&A...376..359H
 CREATOR = 'MODFITS '           / Alexander Menshchikov, DAp IRFU CEA Saclay
 DATE    = '2020-08-05T17:30:38' / creation date and time
 BZERO   =  0.0000000000000E+00 / zero point in scaling equation
 BSCALE  =  1.0000000000000E+00 / linear factor in scaling equation
 DATAMAX =  8.4043195312500E+04 / maximum data value
 DATAMIN = -8.9359838867188E+02 / minimum data value
 BUNIT   = 'MJy/sr  '           / physical units of the array values
 CTYPE1  = 'RA---TAN'           / name of the coordinate axis
 CTYPE2  = 'DEC--TAN'           / name of the coordinate axis
 CRPIX1  =  1.3005000000000E+03 / coordinate system reference pixel
 CRPIX2  =  1.7505000000000E+03 / coordinate system reference pixel
 CROTA1  =  5.1246834126890E-04 / coordinate system rotation angle
 CROTA2  =  5.1246834126890E-04 / coordinate system rotation angle
 CD1_1   = -8.3333333330000E-04 / linear projection matrix
 CD1_2   =  0.0000000000000E+00 / linear projection matrix
 CD2_1   =  0.0000000000000E+00 / linear projection matrix
 CD2_2   =  8.3333333330000E-04 / linear projection matrix
 CRVAL1  =  3.5999968770000E+02 / coordinate value at reference pixel
 CRVAL2  =  3.6225910130000E-04 / coordinate value at reference pixel
 CDELT1  = -8.3333333333333E-04 / coordinate increment along axis
 CDELT2  =  8.3333333333333E-04 / coordinate increment along axis
 OBJECT  = 'Observed region ID' / object identifier
 RA      =  3.5999968770000E+02 / right ascension
 DEC     =  3.6225910130000E-04 / declination
 EQUINOX =  2.0000000000000E+03 / equinox
 WAVE    =  2.5000000000000E+02 / wavelength (microns)
 END
 No more extensions.

 Pixel sizes: -3.000E+00 3.000E+00 arcsec
 
 Done.
\end{verbnobox}
%----------------------------------------------------------------------------------------------------------------------------------
The easiest way of adapting the FITS headers of the new images to the \textsl{getsf} headers is to multiply them by 1:
%----------------------------------------------------------------------------------------------------------------------------------
\begin{verbnobox}[\mbox{}\small]
> modfits multiply 1 universe_250 -o universe_250.m
\end{verbnobox}
%----------------------------------------------------------------------------------------------------------------------------------
A different name (\verb!universe_250.m!) was specified for the output file, in order to keep the original image and header
un\-altered. The keywords \verb!CDELT1!, \verb!CDELT2!, and \verb!BUNIT! are among those that need to be changed most likely. The
supported keywords can be added or modified, one at a time, for example:
%----------------------------------------------------------------------------------------------------------------------------------
\begin{verbnobox}[\mbox{}\small]
> modfits key CDELT1 -2 universe_250.m -o universe_250.m
> modfits key CDELT2 2 universe_250.m -o universe_250.m
\end{verbnobox}
%----------------------------------------------------------------------------------------------------------------------------------
Attention must be paid to the presence of \verb!CDELT1!, \verb!CDELT2!, \verb!CD1_1!, \verb!CD1_2!, \verb!CD2_1!, \verb!CD2_2!, and
\verb!BUNIT!. If the linear projection matrix is absent, then it can be added using \textsl{modfits}, assuming a zero rotation
angle: \verb!CD1_1=CDELT1!, \verb!CD1_2=0!, \verb!CD2_1=0!, \verb!CD2_2=CDELT2!. When the rotation matrix is present, but the
\verb!CDELT1! and \verb!CDELT2! keywords are not found or their values are wrong, they must be specified using \textsl{modfits}.
Only two physical units are supported by \textsl{getsf}: \verb!MJy/sr! and \verb!H2/cm^2!. If the image units are neither of the
two, \textsl{modfits} can be used for their conversion or to add the intensity or surface density units, for example:\\
%----------------------------------------------------------------------------------------------------------------------------------
\vspace{-2mm}\noindent\\
{\small\verb!> modfits Jy/pixel!$\sim$\verb!MJy/sr universe_250.m -o universe_250.m!}\\
{\small\verb!> modfits Jy/beam!$\sim$\verb!MJy/sr 18.2 universe_250.m -o universe_250.m!}\\
{\small\verb!> modfits mJy/beam!$\sim$\verb!MJy/sr 18.2 universe_250.m -o universe_250.m!}\\
{\small\verb!> modfits key BUNIT MJy/sr universe_250.m -o universe_250.m!}\\
{\small\verb!> modfits key BUNIT H2/cm^2 universe_250.m -o universe_250.m!}
%----------------------------------------------------------------------------------------------------------------------------------
\begin{quote} %--------------------------------------------------------------------------------------------------------------------
\textbf{NOTE}. When the observed data set contains just a single image, the reader can skip the next three sections and fast
forward to Sect.~\ref{omasks}.
\end{quote} %----------------------------------------------------------------------------------------------------------------------

%||||||||||||||||||||||||||||||||||||||||||||||||||||||||||||||||||||||||||||||||||||||||||||||||||||||||||||||||||||||||||||||||||

\subsection{Reprojecting images using a reference image}
\label{refimage}

This is the simplest approach to the preparation of the observed images for \textsl{getsf}. All images can be resampled
(reprojected) by \textsl{resample} to the same grid of pixels, using the header information from one of the images, considered as a
reference image. This method can be applied, if the parameters of one of the images (pixel size, numbers of pixels, coordinates,
etc.) are deemed suitable also for all other images. Assuming that the data set contains three images, \verb!universe_250!,
\verb!universe_350!, and \verb!universe_500!, and that one of them (\verb!universe_250!) can be used as a reference image, the
commands look like this:
%----------------------------------------------------------------------------------------------------------------------------------
\begin{verbnobox}[\mbox{}\small]
> resample universe_350 universe_250
> resample universe_500 universe_250
\end{verbnobox}
%----------------------------------------------------------------------------------------------------------------------------------
The resulting images, named \verb!universe_350.resamp! and \verb!universe_500.resamp!, are resampled to the pixel grid of
\verb!universe_250!, with consistent FITS headers. Before using \verb!universe_250! as the reference image, it may be necessary to
optimize (reduce) its size by cutting off image borders that do not contain useful information, for example:
%----------------------------------------------------------------------------------------------------------------------------------
\begin{verbnobox}[\mbox{}\small]
> modfits border -100 universe_250 -o universe_250.m
> modfits border -100,-200,-50,-20 universe_250 -o universe_250.m
\end{verbnobox}
%----------------------------------------------------------------------------------------------------------------------------------
In the first of the two examples, the widths of the borders to cut (in pixels) are the same (\verb!100!) on all sides, whereas the
second example shows that they may be different on each side (left: \verb!100!, right: \verb!200!, bottom: \verb!50!, top:
\verb!20!). For very large images, the \textsl{getsf} computation time (also the storage space) may be an issue, hence it is very
important to minimize the image size, in order to reduce its processing time as much as possible.

%||||||||||||||||||||||||||||||||||||||||||||||||||||||||||||||||||||||||||||||||||||||||||||||||||||||||||||||||||||||||||||||||||

\subsection{Reprojecting images with \textsl{prepobs}}
\label{prepobs}

This approach to the preparation of the observed images for \textsl{getsf} is incorporated in \textsl{prepobs} that internally uses
\textsl{resample} for the actual reprojection of the images. This method is more versatile, with a much larger number of optional
parameters, than when using the reference image (Sect.~\ref{refimage}). It may be recommended to prepare images in their sequence
from the short to long wavelengths. It may be necessary to rename the input files before running \textsl{prepobs}, inserting the
wavelength in the file name (separated by '\verb!_!'), to make sure that the wavelength is understood by the script.

Both plain FITS images and HIPE files with several extensions may be given on input to \textsl{prepobs} (cf. Sect.~\ref{hipecube}).
The common usage is illustrated below, assuming that the data set consists of six \emph{Herschel} images:
%----------------------------------------------------------------------------------------------------------------------------------
\begin{verbnobox}[\mbox{}\small]
> prepobs pacs  20p -scanamorphos Universe universe_070 2 1000 1000 12:35:58.0 +27:57:38.8 0
> prepobs pacs  20  -scanamorphos Universe universe_100 2 1000 1000 12:35:58.0 +27:57:38.8 0
> prepobs pacs  20p -scanamorphos Universe universe_160 2 1000 1000 12:35:58.0 +27:57:38.8 0
> prepobs spire 60p -naive        Universe universe_250 2 1000 1000 12:35:58.0 +27:57:38.8 0
> prepobs spire 60p -naive        Universe universe_350 2 1000 1000 12:35:58.0 +27:57:38.8 0
> prepobs spire 60p -naive        Universe universe_500 2 1000 1000 12:35:58.0 +27:57:38.8 0
\end{verbnobox}
%----------------------------------------------------------------------------------------------------------------------------------
where \verb!pacs! and \verb!spire! refer to the \emph{Herschel} instruments, \verb!20!, \verb!20p! and \verb!60p! are the scanning
speeds of the imaging (20 and 60$\arcsec$s$^{-1}$, \verb!p! indicates the parallel mode), \verb!scanamorphos! and \verb!naive! are
the map-making methods, used to produce the images ('\verb!-!' indicates the plain FITS image), \verb!Universe! is the name of the
observed region, \verb!universe_xxx! are the names of the image files (with the wavelength \verb!xxx! attached), \verb!2! is the
desired pixel size (arcsec), \verb!1000! \verb!1000! are the desired image sizes (pixels), \verb!12:35:58.0! \verb!+27:57:38.8! are
the world coordinates of the field center, and the last \verb!0! is the rotation angle (east of north) of the final prepared images.

The following commands prepare the \textsl{Spitzer} images:
%----------------------------------------------------------------------------------------------------------------------------------
\begin{verbnobox}[\mbox{}\small]
> prepobs irac   1 -other Universe universe_irac1   0.5 1000 1000 12:35:58.0 +27:57:38.8 0
> prepobs irac   2 -other Universe universe_irac2   0.5 1000 1000 12:35:58.0 +27:57:38.8 0
> prepobs irac   3 -other Universe universe_irac3   0.5 1000 1000 12:35:58.0 +27:57:38.8 0
> prepobs irac   4 -other Universe universe_irac4   0.5 1000 1000 12:35:58.0 +27:57:38.8 0
> prepobs mips  24 -other Universe universe_mips24  0.5 1000 1000 12:35:58.0 +27:57:38.8 0
> prepobs mips  70 -other Universe universe_mips70  0.5 1000 1000 12:35:58.0 +27:57:38.8 0
> prepobs mips 160 -other Universe universe_mips160 0.5 1000 1000 12:35:58.0 +27:57:38.8 0
\end{verbnobox}
%----------------------------------------------------------------------------------------------------------------------------------
where \verb!irac! and \verb!mips! refer to the \emph{Spitzer} instruments, \verb!1!, \verb!2!, \verb!3!, \verb!4!, \verb!24!,
\verb!70!, and \verb!160! identifying the bands, \verb!-other! indicates that the map-making method is unknown or unimportant, and
\verb!0.5! is the desired pixel size (arcsec).

The following command prepares the ALMA images in the 1 and 3\,mm bands:
%----------------------------------------------------------------------------------------------------------------------------------
\begin{verbnobox}[\mbox{}\small]
> prepobs other 0.5 -other Universe universe_1mm 0.1 1000 1000 12:35:58.0 +27:57:38.8 0
> prepobs other 0.5 -other Universe universe_3mm 0.1 1000 1000 12:35:58.0 +27:57:38.8 0
\end{verbnobox}
%----------------------------------------------------------------------------------------------------------------------------------
where \verb!other! refers to an instrument (unknown to \textsl{getsf}), \verb!0.5! is the angular resolution (arcsec),
\verb!-other! refers to an map-making method (unknown to \textsl{getsf}), and \verb!0.1! is the desired pixel size (arcsec).

Several runs of \textsl{prepareobs} with different center coordinates, dimensions, and rotation angles are sufficient to optimize
the layout of the observed coverages and to reduce the final size of the prepared images; it is also possible to create a smaller
sub-field from the original large image and to cut useless borders (an alternative to the \textsl{modfits} technique,
Sect.~\ref{refimage}). For very large images, the \textsl{getsf} computation time (also the storage space) may be an issue, hence
it is very important to minimize the image size, in order to reduce its processing time as much as possible.

%||||||||||||||||||||||||||||||||||||||||||||||||||||||||||||||||||||||||||||||||||||||||||||||||||||||||||||||||||||||||||||||||||

\subsection{Verifying alignment of all images}
\label{alignment}
     
In the multi-wavelength extractions with \textsl{getsf}, it is very important to verify alignment of all images. In general, the
prepared data set contains images, observed with different telescopes and/or instruments and created with different map-making
methods. Such differences may easily lead to misalignments of some images. Any substantial shift of the intensity peaks in some
images with respect to the other could lead to the appearance of many spurious sources. Therefore, all prepared images must be
carefully compared with each other, to make sure that relatively bright unresolved peaks in various places across the images stay
at the same pixels. To reveal possible misalignments, it is sufficient to open each pair of the prepared images in \textsl{ds9},
one over the other, and quickly alternate between the two frames, going from the highest-resolution to the lowest-resolution pairs.
If some images are misaligned, such blinking would reveal obvious jumps between the two frames. Misalignments can be corrected by
shifting the problematic images, for example:
%----------------------------------------------------------------------------------------------------------------------------------
\begin{verbnobox}[\mbox{}\small]
> modfits border -1,1,-1,1 universe_350 -o universe_350.shift
\end{verbnobox}
%----------------------------------------------------------------------------------------------------------------------------------
In the above example, the image is shifted by 1 pixel in the diagonal direction to the lower-right. After the attempt to shift one
image, the visual inspection must be repeated, until all images become well aligned.

%||||||||||||||||||||||||||||||||||||||||||||||||||||||||||||||||||||||||||||||||||||||||||||||||||||||||||||||||||||||||||||||||||

\subsection{Creating observational masks}
\label{omasks}

Most astronomical observations have irregularly shaped coverages and limited usable areas that are smaller than the image. To
include in the image processing only the usable areas of good quality, it is necessary to create the observational masks -- special
images with the pixel values of only 1 and 0. Using the masks, \textsl{getsf} processes only the good areas of the original images,
with the mask value of 1. To facilitate the image preparation with \textsl{prepobs} (Sect.~\ref{prepobs}), the script always
creates the simplest default masks (with all pixels of 1). The masks can also be made by running \textsl{modfits}, for example:
%----------------------------------------------------------------------------------------------------------------------------------
\begin{verbnobox}[\mbox{}\small]
> modfits multiply 0 universe_250 -o universe_250.omask
> modfits add 1 universe_250.omask -o universe_250.omask
\end{verbnobox}
%----------------------------------------------------------------------------------------------------------------------------------
Such simple masks may only be used, if \emph{all pixels} of the images are good (usable). However, for most observations the masks
must be prepared manually, for each image independently.

To create the masks manually, it is necessary to use an image manipulation program, such as \textsl{imagej}
\citep{Abramoff_etal2004} or \textsl{gimp} (\url{http://www.gimp.org/}), that allows users to create a polygon over an image,
convert the polygon into a mask, and save it in the FITS format. For example, assuming that \textsl{imagej} is used, here are the
steps: (1) Open one of the prepared images: Menu $\rhd$ File $\rhd$ Open $\rhd$ \verb!universe_250.fits!. (2) Mirror the image
vertically (\textsl{imagej} opens it upside down): Menu $\rhd$ Image $\rhd$ Transform $\rhd$ Flip vertically. (3) Adjust brightness
and contrast: Menu $\rhd$ Image $\rhd$ Adjust $\rhd$ Brightness/contrast. (4) Press the button 'Polygon selections'. (5) Select a
polygon, surrounding the usable area of the image. (6) Convert the polygon in a mask: Menu $\rhd$ Edit $\rhd$ Selection $\rhd$
Create mask. (7) Save the mask: Menu $\rhd$ File $\rhd$ Save as $\rhd$ FITS $\rhd$ \verb!mask_250.fits!.

The resulting mask image, created by \textsl{imagej} has a poor FITS header and inappropriate values (255, instead of 1), therefore
it must be corrected for use with \textsl{getsf}. Assuming that \verb!mask.fits! was saved in the \verb!images! directory with the
prepared \verb!universe_250.fits!, the following two commands correct the mask:
%----------------------------------------------------------------------------------------------------------------------------------
\begin{verbnobox}[\mbox{}\small]
> operate universe_250 h mask_250 -o universe_250.omask
> modfits divide 255 universe_250.omask -o universe_250.omask
\end{verbnobox}
%----------------------------------------------------------------------------------------------------------------------------------
As the result, the following pair of files (with identical headers) must exist in the \verb!images! directory:
%----------------------------------------------------------------------------------------------------------------------------------
\begin{verbnobox}[\mbox{}\small]
universe_250.fits
universe_250.omask.fits
\end{verbnobox}
%----------------------------------------------------------------------------------------------------------------------------------
The same steps must be repeated for all prepared images. Consistency of each pair must be verified by opening the prepared image in
together with its observational mask in \textsl{ds9} and blinking them to make sure that the mask indeed covers only the good area
of the images, for which the user plans to run a \textsl{getsf} extraction. It is also necessary to check that the images do not
have wide empty areas around the border, with zero values in the observational masks. If such empty borders exist, it is strongly
recommended to reduce the image size by cutting the borders, for example:
%----------------------------------------------------------------------------------------------------------------------------------
\begin{verbnobox}[\mbox{}\small]
> modfits border -50,-30,-40,-60 universe_250 -o universe_250.m
> modfits border -50,-30,-40,-60 universe_250.omask -o universe_250.omask.m
\end{verbnobox}
%----------------------------------------------------------------------------------------------------------------------------------
It is important to emphasize that the observational masks must have topologically simple shapes, without any isolated 'holes' (zero
areas) inside the masked field. All zero pixels of the mask image must be connected to the image edges.

%||||||||||||||||||||||||||||||||||||||||||||||||||||||||||||||||||||||||||||||||||||||||||||||||||||||||||||||||||||||||||||||||||
%||||||||||||||||||||||||||||||||||||||||||||||||||||||||||||||||||||||||||||||||||||||||||||||||||||||||||||||||||||||||||||||||||

\section{Deriving high-resolution surface densities}
\label{hires}

\begin{quote} %--------------------------------------------------------------------------------------------------------------------
\textbf{NOTE}. If the prepared data set contains only one or two images or it does not allow derivation of surface densities (by
pixel fitting) for other reasons, the reader can skip this section and fast forward to Sect.~\ref{editconfig}.
\end{quote} %----------------------------------------------------------------------------------------------------------------------
A high-resolution surface density image can be derived, using \textsl{hires}, from the multi-wavelength far-infrared and
submillimeter continuum images, such as those obtained in the \emph{Herschel} observations. The density image may be as sharp as
the shortest-wavelength image of a sufficient quality, with the angular resolution reaching $5.6{\arcsec}$, when using the
\emph{Herschel} $70$\,$\mu$m image. If the image is too noisy or unusable for other reasons, the highest resolution would be
defined by the $100$ or $160$\,$\mu$m image ($6.8{-}11.3{\arcsec}$). In the multi-wavelength extractions with \textsl{getsf}, it
may be recommended to use the (optional) high-resolution surface density for better detections and deblending of dense structures.
It would be advantageous to have it complement the original data set, as an additional image, ``observed'' in a fictitious
waveband. For more details of the algorithm, refer to \cite{Men'shchikov2021getsf}.

%||||||||||||||||||||||||||||||||||||||||||||||||||||||||||||||||||||||||||||||||||||||||||||||||||||||||||||||||||||||||||||||||||

\subsection{The \textsl{hires} configuration file}
\label{hiconfig}

To prepare for the derivation of surface density images, it is necessary to create the default configuration file by executing
\textsl{hires} in a dedicated subdirectory, preferably within the \verb!images! directory:
%----------------------------------------------------------------------------------------------------------------------------------
\begin{verbnobox}[\mbox{}\small]
> mkdir -p hires
> cd hires
> hires +hires.cfg
\end{verbnobox}
%----------------------------------------------------------------------------------------------------------------------------------
This command creates the configuration file named \verb!+hires.cfg!. When opening files with '\verb!+!' in their names in the text
editors interpreting the '\verb!+!' character as a special symbol, it is necessary to append the '\verb!./!' string to the file
name, to tell the editor that the character is a part of the file name, for example:
%----------------------------------------------------------------------------------------------------------------------------------
\begin{verbnobox}[\mbox{}\small]
> nedit ./+hires.cfg
\end{verbnobox}
%----------------------------------------------------------------------------------------------------------------------------------
The contents of the default configuration file:
%----------------------------------------------------------------------------------------------------------------------------------
\begin{verbnobox}[\mbox{}\small]
#
# HIRES: CONFIGURATION PARAMETERS
#
# PM  WAVE  BEAM  OFFSET  IMAGE      MASKIMAGE
# +/- (um)   (")  MJy/sr  file       file
#
   -   070   8.4   0.00   image_070  image_070.omask
   -   100   9.4   0.00   image_100  image_100.omask
   +   160  13.5   0.00   image_160  image_160.omask
   +   250  18.2   0.00   image_250  image_250.omask
   +   350  24.9   0.00   image_350  image_350.omask
   +   500  36.3   0.00   image_500  image_500.omask
\end{verbnobox}
%----------------------------------------------------------------------------------------------------------------------------------
The files with images and observational masks are those created in the \verb!images! directory (Sect.~\ref{preparation}), specified
with either relative or absolute path. If some of the wavebands are not present in the set of prepared images (e.g., the
100\,${\mu}$m image), the corresponding lines must be deleted.

%||||||||||||||||||||||||||||||||||||||||||||||||||||||||||||||||||||||||||||||||||||||||||||||||||||||||||||||||||||||||||||||||||

\subsection{The configuration parameters}
\label{edithiconf}

The meaning of the first line of parameters in \verb!+hires.cfg! is explained below.
%----------------------------------------------------------------------------------------------------------------------------------
\begin{verbnobox}[\mbox{}\small]
   -   070   8.4   0.00   image_070  image_070.omask
\end{verbnobox}
%----------------------------------------------------------------------------------------------------------------------------------
The 1st parameter (`\verb!-!') indicates that the 70\,${\mu}$m image will not be used for the fitting of the spectral shape of each
pixel. Otherwise, the short-wavelength image would lead to the overestimated temperatures, hence underestimated surface densities.
In general, images at the wavelengths shorter than 160\,${\mu}$m should not be used for the fitting. The 2nd parameter (\verb!070!)
is the wavelength of the corresponding image, with a leading zero. The 3rd parameter (\verb!8.4!) is the angular resolution, i.e.,
the observational beam size (FWHM, in arcsecs). The 4th parameter (\verb!0.00!) is the offset (in MJy/sr) for the image, correcting
the image absolute calibration. In general, the default zero must be replaced with an estimated non-zero value, which is likely to
be a separate non-trivial problem. However, the value may be critically important for the acceptably good quality of the resulting
surface density image. The 5th parameter (\verb!image_070!) is the prepared 70\,${\mu}$m image (\verb!images/universe_070!) and the
6th parameter its corresponding observational mask (Sect.~\ref{preparation}). When prepared, the configuration file may look like:
%----------------------------------------------------------------------------------------------------------------------------------
\begin{verbnobox}[\mbox{}\small]
#
# HIRES: CONFIGURATION PARAMETERS
#
# PM  WAVE  BEAM  OFFSET  IMAGE               MASKIMAGE
# +/- (um)   (")  MJy/sr  file                file
#
   -   070   8.4   400   images/universe_070  images/universe_070.omask
   +   160  13.5   850   images/universe_160  images/universe_160.omask
   +   250  18.2     0   images/universe_250  images/universe_250.omask
   +   350  24.9    50   images/universe_350  images/universe_350.omask
   +   500  36.3   160   images/universe_500  images/universe_500.omask
\end{verbnobox}
%----------------------------------------------------------------------------------------------------------------------------------

%||||||||||||||||||||||||||||||||||||||||||||||||||||||||||||||||||||||||||||||||||||||||||||||||||||||||||||||||||||||||||||||||||

\subsection{Convolving images to lower resolutions}
\label{convolving}

The following command instructs \verb!hires! to convolve all input images to all lower angular resolutions found in
\verb!+hires.cfg! and to create the configuration files \verb!+hires.r250.cfg!, \verb!+hires.r350.cfg!, and \verb!+hires.r500.cfg!,
necessary for the next step of fitting the spectral shapes of images pixels:
%----------------------------------------------------------------------------------------------------------------------------------
\begin{verbnobox}[\mbox{}\small]
> hires +hires.cfg prepare
\end{verbnobox}
%----------------------------------------------------------------------------------------------------------------------------------
As a result, the following set of convolved images is created in the \verb!hires! directory:\\
%----------------------------------------------------------------------------------------------------------------------------------
\vspace{-2mm}\noindent\\
{\small\verb!universe_070.r070!} $\,\,\leadsto$ 070 $\mu$m image at the resolution of the 070 $\mu$m image\\
{\small\verb!universe_070.r160!} $\,\,\leadsto$ 070 $\mu$m image at the resolution of the 160 $\mu$m image\\
{\small\verb!universe_070.r250!} $\,\,\leadsto$ 070 $\mu$m image at the resolution of the 250 $\mu$m image\\
{\small\verb!universe_070.r350!} $\,\,\leadsto$ 070 $\mu$m image at the resolution of the 350 $\mu$m image\\
{\small\verb!universe_070.r500!} $\,\,\leadsto$ 070 $\mu$m image at the resolution of the 500 $\mu$m image\\
{\small\verb!universe_160.r160!} $\,\,\leadsto$ 160 $\mu$m image at the resolution of the 160 $\mu$m image\\
{\small\verb!universe_160.r250!} $\,\,\leadsto$ 160 $\mu$m image at the resolution of the 250 $\mu$m image\\
{\small\verb!universe_160.r350!} $\,\,\leadsto$ 160 $\mu$m image at the resolution of the 350 $\mu$m image\\
{\small\verb!universe_160.r500!} $\,\,\leadsto$ 160 $\mu$m image at the resolution of the 500 $\mu$m image\\
{\small\verb!universe_250.r250!} $\,\,\leadsto$ 250 $\mu$m image at the resolution of the 250 $\mu$m image\\
{\small\verb!universe_250.r350!} $\,\,\leadsto$ 250 $\mu$m image at the resolution of the 350 $\mu$m image\\
{\small\verb!universe_250.r500!} $\,\,\leadsto$ 250 $\mu$m image at the resolution of the 500 $\mu$m image\\
{\small\verb!universe_350.r350!} $\,\,\leadsto$ 350 $\mu$m image at the resolution of the 350 $\mu$m image\\
{\small\verb!universe_350.r500!} $\,\,\leadsto$ 350 $\mu$m image at the resolution of the 500 $\mu$m image\\
{\small\verb!universe_500.r500!} $\,\,\leadsto$ 500 $\mu$m image at the resolution of the 500 $\mu$m image
%----------------------------------------------------------------------------------------------------------------------------------

%||||||||||||||||||||||||||||||||||||||||||||||||||||||||||||||||||||||||||||||||||||||||||||||||||||||||||||||||||||||||||||||||||

\subsection{Fitting spectral shapes of pixels with \textsl{fitfluxes}}
\label{pixfitting}

The fitting method used by \textsl{fitfluxes} is described in \cite{Men'shchikov2016}; its usage information is shown, when the
utility is run without parameters. The following command lists all \textsl{fitfluxes} runs and parameters, necessary to produce all
variants of the surface density and temperature images with different resolutions:
%----------------------------------------------------------------------------------------------------------------------------------
\begin{verbnobox}[\mbox{}\small]
> hires +hires.cfg fitfluxes
\end{verbnobox}
%----------------------------------------------------------------------------------------------------------------------------------
As a result, the user screen displays the following information on the parameters to supply to \textsl{fitfluxes} and on the
correct sequence of the commands:
%----------------------------------------------------------------------------------------------------------------------------------
\begin{verbnobox}[\mbox{}\small]
 Here are the commands fitting the spectral shapes of pixels in observed images
 to compute the high-resolution surface density and temperature images using the
 HIRES algorithm, described by Men’shchikov (2021) A&A, 649, A89:
 
 FITTING PIXEL SPECTRAL SHAPES IN IMAGES AT 500 µm RESOLUTION:

 fitfluxes 0 0 0 -1:070,-2:160,-3:250,4:350,5:500 1:0.2,1:0.2,1:0.1,1:0.1,1:0.1 \
           thinbody 0:1,0:1,2:0 1 10 0.01 +hires.r500.cfg -verb0
 fitfluxes 0 0 0 -1:070,-2:160,3:250,4:350,5:500 1:0.2,1:0.2,1:0.1,1:0.1,1:0.1 \
           thinbody 0:1,0:1,2:0 1 10 0.01 +hires.r500.cfg -verb0
 fitfluxes 0 0 0 -1:070,2:160,3:250,4:350,5:500 1:0.2,1:0.2,1:0.1,1:0.1,1:0.1 \
           thinbody 0:1,0:1,2:0 1 10 0.01 +hires.r500.cfg -verb0

 FITTING PIXEL SPECTRAL SHAPES IN IMAGES AT 350 µm RESOLUTION:

 fitfluxes 0 0 0 -1:070,-2:160,3:250,4:350 1:0.2,1:0.2,1:0.1,1:0.1 \
           thinbody 0:1,0:1,2:0 1 10 0.01 +hires.r350.cfg -verb0
 fitfluxes 0 0 0 -1:070,2:160,3:250,4:350 1:0.2,1:0.2,1:0.1,1:0.1 \
           thinbody 0:1,0:1,2:0 1 10 0.01 +hires.r350.cfg -verb0

 FITTING PIXEL SPECTRAL SHAPES IN IMAGES AT 250 µm RESOLUTION:

 fitfluxes 0 0 0 -1:070,2:160,3:250 1:0.2,1:0.2,1:0.1 \
           thinbody 0:1,0:1,2:0 1 10 0.01 +hires.r250.cfg -verb0

 FITTING PIXEL SPECTRAL SHAPES IN IMAGES AT 500 µm RESOLUTION (OFFSETS CONSISTENCY CHECKS):

 fitfluxes 0 0 0 -1:160,2:250,3:350,-4:500 1:0.2,1:0.1,1:0.1,1:0.1 \
           thinbody 0:1,0:1,2:0 1 10 0.01 +hires.r500.cfg -verb0
 fitfluxes 0 0 0 1:160,2:250,-3:350,-4:500 1:0.2,1:0.1,1:0.1,1:0.1 \
           thinbody 0:1,0:1,2:0 1 10 0.01 +hires.r500.cfg -verb0

 Start the commands (in the sequence shown above) on different processors, in parallel.
\end{verbnobox}
%----------------------------------------------------------------------------------------------------------------------------------
The \textsl{fitfluxes} runs are not executed by \textsl{hires}, because it is much more efficient for the user to run the jobs in
parallel in different terminal sessions (on different processors), to greatly accelerate the processing. The parallel jobs must be
started \emph{in the sequence shown, one after the other}, because the numbering of the output files depends on that. 
As a result, the following sets of the surface densities and temperatures are created in
the \verb!hires! directory:\\
%----------------------------------------------------------------------------------------------------------------------------------
\vspace{-2mm}\noindent\\
{\small\verb!thin.+hires.rWWW.cfg.0N.chisqnu.fits!} $\,\,\leadsto$ $\chi^2$ values for the resulting fit\\
{\small\verb!thin.+hires.rWWW.cfg.0N.log!} $\,\,\leadsto$ log file for the run\\
{\small\verb!thin.+hires.rWWW.cfg.0N.surfden.fits!} $\,\,\leadsto$ surface density at the resolution of the \verb!WWW! $\mu$m image\\
{\small\verb!thin.+hires.rWWW.cfg.0N.surferr.fits!} $\,\,\leadsto$ relative errors in surface densities\\
{\small\verb!thin.+hires.rWWW.cfg.0N.temperr.fits!} $\,\,\leadsto$ relative errors in temperatures\\
{\small\verb!thin.+hires.rWWW.cfg.0N.tempers.fits!\vspace{-2mm}} $\,\,\leadsto$ temperature at the resolution of the \verb!WWW! $\mu$m image\\\\
%----------------------------------------------------------------------------------------------------------------------------------
In the above names, \verb!WWW! is the wavelength and \verb!0N! is the sequential number (\verb!00!, \verb!01!, \verb!02!,
\verb!03!, \verb!04!) of the \textsl{fitfluxes} run. When the jobs are cancelled before they have finished, all log files with the
names \verb!*.cfg.??.log! must be deleted and the sequence of the \textsl{fitfluxes} jobs must be started again.

%||||||||||||||||||||||||||||||||||||||||||||||||||||||||||||||||||||||||||||||||||||||||||||||||||||||||||||||||||||||||||||||||||

\subsection{Computing high-resolution surface densities}
\label{derivation}

The derived set of surface densities with three angular resolutions is used to differentially increase their resolution to any
value among the resolutions present in \verb!+hires.cfg!:
%----------------------------------------------------------------------------------------------------------------------------------
\begin{verbnobox}[\mbox{}\small]
> hires +hires.cfg complete
\end{verbnobox}
%----------------------------------------------------------------------------------------------------------------------------------
The command creates the high-resolution surface densities and temperatures (in \verb!+hires/+results!), employing the algorithm
published by \cite{Men'shchikov2021getsf}, a substantial improvement over the previous algorithm \citep{Palmeirim_etal2013}:\\
%----------------------------------------------------------------------------------------------------------------------------------
\vspace{-2mm}\noindent\\
{\small\verb!hi.c.surface.density.r8p4.fits !} $\,\,\leadsto$ high-resolution high-contrast density at 8.4$\arcsec$ resolution\\
{\small\verb!hi.c.surface.density.r9p4.fits !} $\,\,\leadsto$ high-resolution high-contrast density at 9.4$\arcsec$ resolution\\
{\small\verb!hi.c.surface.density.r13p5.fits!} $\,\,\leadsto$ high-resolution high-contrast density at 13.5$\arcsec$ resolution\\
{\small\verb!hi.c.surface.density.r18p2.fits!} $\,\,\leadsto$ high-resolution high-contrast density at 18.2$\arcsec$ resolution\\
{\small\verb!hi.c.surface.density.r24p9.fits!} $\,\,\leadsto$ high-resolution high-contrast density at 24.9$\arcsec$ resolution\\
{\small\verb!hi.c.surface.density.r36p3.fits!} $\,\,\leadsto$ high-resolution high-contrast density at 36.3$\arcsec$ resolution\\
{\small\verb!hi.surface.density.r8p4.fits !} $\,\,\leadsto$ high-resolution density at 8.4$\arcsec$ resolution\\
{\small\verb!hi.surface.density.r9p4.fits !} $\,\,\leadsto$ high-resolution density at 9.4$\arcsec$ resolution\\
{\small\verb!hi.surface.density.r13p5.fits!} $\,\,\leadsto$ high-resolution density at 13.5$\arcsec$ resolution\\
{\small\verb!hi.surface.density.r18p2.fits!} $\,\,\leadsto$ high-resolution density at 18.2$\arcsec$ resolution\\
{\small\verb!hi.surface.density.r24p9.fits!} $\,\,\leadsto$ high-resolution density at 24.9$\arcsec$ resolution\\
{\small\verb!hi.surface.density.r36p3.fits!\vspace{-2mm}} $\,\,\leadsto$ high-resolution density at 36.3$\arcsec$ resolution\\\\
%----------------------------------------------------------------------------------------------------------------------------------
The high-contrast variants of the high-resolution images (\verb!hi.c.surface.density.*.fits!) with an artificially enhanced
contrast may be useful for the analysis of the fine unresolved structures, present in the images, but they may not be suitable for
the \textsl{getsf} extractions. It is recommended to use only the regular versions of the images (\verb!hi.surface.density.*.fits!)
in the source and filament extractions with \textsl{getsf}. 

\begin{quote} %--------------------------------------------------------------------------------------------------------------------
\textbf{NOTE}. When choosing the best angular resolution of the images, the user should verify their quality very carefully. This
is especially critical for the images with the highest resolutions of 8.4$\arcsec$ and 9.4$\arcsec$, whose quality may be
unacceptably low due to excessively high noise or artifacts.
\end{quote} %----------------------------------------------------------------------------------------------------------------------

%||||||||||||||||||||||||||||||||||||||||||||||||||||||||||||||||||||||||||||||||||||||||||||||||||||||||||||||||||||||||||||||||||

\subsection{Observational mask for surface densities}
\label{densobsmask}

The \verb!+hires/+results! contains an observational mask, automatically created for only the highest-resolution surface density
image, i.e., for \verb!hi.surface.density.r8p4!. If another image is planned for an extraction (e.g.,
\verb!hi.surface.density.r13p5! or \verb!hi.surface.density.r18p2!), then the mask must either be copied or renamed, for example:
%----------------------------------------------------------------------------------------------------------------------------------
\begin{verbnobox}[\mbox{}\small]
> cp +results/hi.surface.density.r8p4.omask.fits +results/hi.surface.density.r13p5.omask.fits
> mv +results/hi.surface.density.r8p4.omask.fits +results/hi.surface.density.r18p2.omask.fits
\end{verbnobox}
%----------------------------------------------------------------------------------------------------------------------------------
The surface density mask is created as the intersection of the masks of the individual images (Sect.~\ref{omasks}) specified in the
configuration file \verb!+hires.cfg! (Sect.~\ref{hiconfig}), therefore it does not require any adjustments.

%||||||||||||||||||||||||||||||||||||||||||||||||||||||||||||||||||||||||||||||||||||||||||||||||||||||||||||||||||||||||||||||||||
%||||||||||||||||||||||||||||||||||||||||||||||||||||||||||||||||||||||||||||||||||||||||||||||||||||||||||||||||||||||||||||||||||

\section{Configuring source and filament extraction}
\label{editconfig}

In order to complete the preparation for the \textsl{getsf} extraction, it is necessary to create the default configuration file
\verb!+getsf.cfg! by executing \textsl{getsf} in a directory that does not contain that file:
%----------------------------------------------------------------------------------------------------------------------------------
\begin{verbnobox}[\mbox{}\small]
> cd ../241122/runs
> getsf
\end{verbnobox}
%----------------------------------------------------------------------------------------------------------------------------------

%||||||||||||||||||||||||||||||||||||||||||||||||||||||||||||||||||||||||||||||||||||||||||||||||||||||||||||||||||||||||||||||||||

\subsection{The \textsl{getsf} configuration file}
\label{config}

The content of the default configuration file \verb!+getsf.cfg! facilitates preparations for the 6 \emph{Herschel} wavebands.
However, the information can be easily modified for use with any other set of images. When opening files with '\verb!+!' in their
names in the text editors interpreting the '\verb!+!' character as a special symbol, it is necessary to append the '\verb!./!'
string to the file name, to tell the editor that the character is a part of the file name, for example:
%----------------------------------------------------------------------------------------------------------------------------------
\begin{verbnobox}[\mbox{}\small]
> nedit ./+getsf.cfg
\end{verbnobox}
%----------------------------------------------------------------------------------------------------------------------------------
The contents of the default configuration file:
\begin{verbnobox}[\mbox{}\small]
#_______________________________________________________________________________________________
#                                                                                               
# GETSF • Multi-Scale Multi-Wavelength Source & Filament Extraction • Version 241122 • Alexander
#_______________________________________________________________________________________________
#                    |                                      !
# VALUES             | PARAMETERS                           ! COMMENT  NB: default beams are for
#____________________|______________________________________!___________________________________
#                    |                                      !
 fieldname 100       | prefix distance .................... ! prefix for naming files (never use
 6                   | nwmax ...............{ nwmax < 100 } ! maximum number of wavelengths spec
 6 1 2 3 4 5 6       | nwaves nw[i] ....................... ! number of waveblengths and their i
 070  8.4 s070 f070 n col | ............................... ! wave, beam, [-]srcmaxsize, [-]film
 /path/to/image ../070 | .................................. ! paths to image and monochromatic d
 image_070           | .................................... ! prepared image name
 100  9.4 s100 f100 n col | ............................... ! wave, beam, [-]srcmaxsize, [-]film
 /path/to/image ../100 | .................................. ! paths to image and monochromatic d
 image_100           | .................................... ! prepared image name
 160 13.5 s160 f160 n col | ............................... ! wave, beam, [-]srcmaxsize, [-]film
 /path/to/image ../160 | .................................. ! paths to image and monochromatic d
 image_160           | .................................... ! prepared image name
 250 18.2 s250 f250 n col | ............................... ! wave, beam, [-]srcmaxsize, [-]film
 /path/to/image ../250 | .................................. ! paths to image and monochromatic d
 image_250           | .................................... ! prepared image name
 350 24.9 s350 f350 n col | ............................... ! wave, beam, [-]srcmaxsize, [-]film
 /path/to/image ../350 | .................................. ! paths to image and monochromatic d
 image_350           | .................................... ! prepared image name
 500 36.3 s500 f500 n col | ............................... ! wave, beam, [-]srcmaxsize, [-]film
 /path/to/image ../500 | .................................. ! paths to image and monochromatic d
 image_500           | .................................... ! prepared image name
 y y                 | sources filaments .................. ! extract sources and/or filaments {
 y 1 1 50            | separate its1 itf1 niter ........... ! separate components and background
 y 1 1               | flatten itsd1 itfd1 ................ ! flatten background-subtracted comp
 y                   | complete ........................... ! complete separation and flattening
 y                   | combine ............................ ! combine decomposed structures over
 y 2                 | detect skelsig ..................... ! detect structures; skeleton signif
 y 1 1 50 n          | measure goodmin itn itx sproima .... ! measure structures; minimum goodne
 y                   | visualize .......................... ! visualize sources
 5                   | nsigmacombos (!expert) ............. ! number of sigmas for combining sou
 2                   | nsigmacombof (!expert) ............. ! number of sigmas for combining fil
 1                   | nwavesdetect ....................... ! required minimum number of wavelen
 2 ?                 | savespace delete ................... ! space saving level {0|1|2}: clean 
 0                   | verbosity .......................... ! verbosity level {0|1|2} for the sc
#_______________________________________________________________________________________________
\end{verbnobox}
%----------------------------------------------------------------------------------------------------------------------------------

%||||||||||||||||||||||||||||||||||||||||||||||||||||||||||||||||||||||||||||||||||||||||||||||||||||||||||||||||||||||||||||||||||

\subsection{The configuration parameters}
\label{params}

The meaning of each line of parameters in \verb!+getsf.cfg! is explained below.
%----------------------------------------------------------------------------------------------------------------------------------
\begin{verbnobox}[\mbox{}\small]
 fieldname 100       | prefix distance .................... !
\end{verbnobox}
%----------------------------------------------------------------------------------------------------------------------------------
The extraction must be given a name that will serve as the prefix for naming the output files. It makes sense to replace the
placeholder \verb!fieldname! with the name of the observed region, for example, \verb!universe!. The number \verb!100! is the
distance (in parsecs) to the region, used to compute the masses and linear densities of filaments. If the distance is unknown or
unimportant, the default value may be kept.
%----------------------------------------------------------------------------------------------------------------------------------
\begin{verbnobox}[\mbox{}\small]
 6                   | nwmax ...............{ nwmax < 100 } !
\end{verbnobox}
%----------------------------------------------------------------------------------------------------------------------------------
This number (\verb!nwmax!) is the total number of images (wavebands) specified in this configuration file. This value is suitable
for the default set of 6 \emph{Herschel} images (the upper limit is 99).
%----------------------------------------------------------------------------------------------------------------------------------
\begin{verbnobox}[\mbox{}\small]
 6 1 2 3 4 5 6       | nwaves nw[i] ....................... !
\end{verbnobox}
%----------------------------------------------------------------------------------------------------------------------------------
The 1st number (\verb!nwaves!) is the maximum number of images (wavebands) to be used in the current run of \textsl{getsf}. The
other six numbers refer to the corresponding images used in the current run. This default line of parameters means that the
extraction will be done sequentially, processing the next image only after the previous one has been completed
(Sect.~\ref{sequence}). This is the easiest way of running \textsl{getsf}, but it makes sense for only relatively small images. For
large images, it is recommended to always process the images in parallel independent \textsl{getsf} runs (Sect.~\ref{parallel}) to
strongly reduce the total computation time, approximately by the factor \verb!nwmax!. Hence, for the six parallel jobs, this line
of parameters takes the following values:\\
%----------------------------------------------------------------------------------------------------------------------------------
\vspace{-2mm}\noindent\\
{\small\verb! 1 1                 | nwaves nw[i] ....................... ! !} $\,\,\leadsto$ for the 1st parallel run\\
{\small\verb! 1 2                 | nwaves nw[i] ....................... ! !} $\,\,\leadsto$ for the 2nd parallel run\\
{\small\verb! 1 3                 | nwaves nw[i] ....................... ! !} $\,\,\leadsto$ for the 3rd parallel run\\
{\small\verb! 1 4                 | nwaves nw[i] ....................... ! !} $\,\,\leadsto$ for the 4th parallel run\\
{\small\verb! 1 5                 | nwaves nw[i] ....................... ! !} $\,\,\leadsto$ for the 5th parallel run\\
{\small\verb! 1 6                 | nwaves nw[i] ....................... ! !\vspace{-2mm}} $\,\,\leadsto$ for the 6th parallel run\\\\
%----------------------------------------------------------------------------------------------------------------------------------
More information on the source and filament extraction with \textsl{getsf} is found in Sect.~\ref{extract}.
%----------------------------------------------------------------------------------------------------------------------------------
\begin{verbnobox}[\mbox{}\small]
 070  8.4 s070 f070 n col | ............................... !
\end{verbnobox}
%----------------------------------------------------------------------------------------------------------------------------------
The 1st number (\verb!070!) is the wavelength (in microns) of the corresponding image with a leading zero. The number of digits
must be the same for all wavelengths, hence the leading zero. The 2nd number (\verb!8.4!) is the angular resolution, i.e., the
observational beam size (FWHM, in arcsecs). All default beams correspond to the \emph{Herschel} fast parallel (PACS/SPIRE) scanning
mode (60$\arcsec$s$^{-1}$). The 3rd parameter (\verb!s070!) is the placeholder for the maximum size (FWHM, arcsec) of the sources
to extract in this image. The 4th parameter (\verb!f070!) is the placeholder for the maximum width (FWHM, arcsec) of the filaments
to extract in this image. If a user is interested in only unresolved sources and filaments, both \verb!s070! and \verb!f070! must
be set to the angular resolution. If there are no filaments or there is no interest in filaments, then \verb!f070! should be set to
the angular resolution \emph{and} processing of filaments must be disabled by setting \verb!filaments! to `\verb!n!' (see below).
For more information on the maximum size (width) of sources (filaments), refer to Sect.~3.1.3 in \cite{Men'shchikov2021getsf}. The
maximum size may also be entered as negative values to stipulate that the image must not be used for the detection of sources
(\verb!-s070!) or filaments (\verb!-f070!). Then the image will only be used for the measurements of the sources or filaments,
detected in the other images. The 5th parameter (`\verb!n!') disables the \emph{Herschel} aperture corrections, therefore it must
be enabled (changed to `\verb!y!') for the \emph{Herschel} images. The last \emph{optional} parameter (\verb!col!) is the
placeholder for the desired color of sources and skeletons in the \textsl{ds9} region files \verb!*.reg!, created by \textsl{getsf}
to aid users in visualizations and analyses. The \textsl{getsf} colors are abbreviated as `\verb!gre!', `\verb!blu!', `\verb!ora!',
`\verb!pur!', `\verb!cya!', `\verb!mag!', `\verb!lim!', `\verb!pin!', `\verb!tea!', `\verb!lav!', `\verb!bro!', `\verb!bei!',
`\verb!mar!', `\verb!min!', `\verb!oli!', `\verb!apr!', `\verb!gra!', `\verb!red!', `\verb!yel!', `\verb!whi!', and `\verb!bla!'.
If the optional parameter is not specified, then the default colors are used.
%----------------------------------------------------------------------------------------------------------------------------------
\begin{verbnobox}[\mbox{}\small]
 /path/to/image ../070 | .................................. !
\end{verbnobox}
%----------------------------------------------------------------------------------------------------------------------------------
The 1st parameter (\verb!/path/to/image!) is the absolute or relative path to the prepared image, i.e., to the \verb!images!
directory (Sect.~\ref{preparation}). The 2nd parameter (\verb!../070!) is the relative path to the monochromatic directory, created
by \textsl{getsf}, where all files related to the processing of the image are created. The two paths must be separated by space,
the monochromatic directory must be separated by space from the vertical line (\verb!|!) that follows it, and \emph{the trailing
slash} (\verb!/!) \emph{in the paths must never be used}.
%----------------------------------------------------------------------------------------------------------------------------------
\begin{verbnobox}[\mbox{}\small]
 image_070           | .................................... !
\end{verbnobox}
%----------------------------------------------------------------------------------------------------------------------------------
The name of the observed image, without the \verb!.fits! extension, that was prepared in the \verb!images! directory, i.e.,
\verb!universe_070!. The observational mask is not specified, because it has the same name as the image, with \verb!.omask!
attached (Sect.~\ref{omasks}).
%----------------------------------------------------------------------------------------------------------------------------------
\begin{verbnobox}[\mbox{}\small]
 100  9.4 s100 f100 n col | ............................... !
 /path/to/image ../100 | .................................. !
 image_100           | .................................... !
 160 13.5 s160 f160 n col | ............................... !
 /path/to/image ../160 | .................................. !
 image_160           | .................................... !
 250 18.2 s250 f250 n col | ............................... !
 /path/to/image ../250 | .................................. !
 image_250           | .................................... !
 350 24.9 s350 f350 n col | ............................... !
 /path/to/image ../350 | .................................. !
 image_350           | .................................... !
 500 36.3 s500 f500 n col | ............................... !
 /path/to/image ../500 | .................................. !
 image_500           | .................................... !
\end{verbnobox}
%----------------------------------------------------------------------------------------------------------------------------------
All other wavebands are edited the same way as the first one. If \verb!nwmax! is smaller than 6, the unused entries for the extra
wavebands of the default configuration file must be deleted.
%----------------------------------------------------------------------------------------------------------------------------------
\begin{verbnobox}[\mbox{}\small]
 y y                 | sources filaments .................. !
\end{verbnobox}
%----------------------------------------------------------------------------------------------------------------------------------
The 1st parameter (\verb!sources!) enables the extraction of sources and the 2nd parameter (\verb!filaments!) enables the
extraction of filaments. Changing either of them to \verb!n! disables the processing of the respective structural component in
\textsl{getsf}. If the user is interested only in the source extraction, the 2nd parameter may be changed to `\verb!n!'. However,
if the user is interested only in the filament extraction, it is not recommended to disable the source extraction, because the
extracted filaments would remain contaminated by the sources. Therefore, if the user is interested in the filament extraction, both
parameters should remain `\verb!y!'.
%----------------------------------------------------------------------------------------------------------------------------------
\begin{verbnobox}[\mbox{}\small]
 y 1 1 50            | separate its1 itf1 niter ........... !
\end{verbnobox}
%----------------------------------------------------------------------------------------------------------------------------------
The 1st parameter (\verb!separate!) enables the separation step of the structural components. This parameter may be changed to
`\verb!n!', when \textsl{getsf} must be restarted to process only the other steps. When the separation step is enabled, but the
image processing must be restarted for any reason, the 2nd parameter (\verb!its1!) defines the number of the iteration to restart
with for the component of sources and the 3rd parameter (\verb!itf1!) defines the same for the component of filaments. When
\verb!its1! or \verb!itf1! is set to 0, then the corresponding iterations are bypassed. The last parameter (\verb!niter!) defines
the maximum number of iterations (\verb!50!) in the separation of the components. If the user is interested only in the filament
extraction, it is possible to reduce the \textsl{getsf} execution time by limiting the source processing only to this separation
step.
%----------------------------------------------------------------------------------------------------------------------------------
\begin{verbnobox}[\mbox{}\small]
 y 1 1               | flatten itsd1 itfd1 ................ !
\end{verbnobox}
%----------------------------------------------------------------------------------------------------------------------------------
The 1st parameter (\verb!flatten!) enables the flattening step of the structural components; this parameter may be changed to
`\verb!n!', when \textsl{getsf} must be restarted to process only the other steps. When the flattening step is enabled, but the
image processing must be restarted for any reason, the 2nd parameter (\verb!itsd1!) defines the number of the iteration to restart
with for the component of sources and the 3rd parameter (\verb!itfd1!) defines the same for the component of filaments. When
\verb!its1! or \verb!itf1! is set to 0, then the corresponding iterations are bypassed. The maximum number of iterations in the
flattening of the components is defined by \verb!niter! from the separation step.
%----------------------------------------------------------------------------------------------------------------------------------
\begin{verbnobox}[\mbox{}\small]
 y                   | complete ........................... !
\end{verbnobox}
%----------------------------------------------------------------------------------------------------------------------------------
This parameter (\verb!complete!) enables the completion step after the separation and flattening of the structural components; this
parameter may be changed to `\verb!n!', when \textsl{getsf} must be restarted to process only the other steps.
%----------------------------------------------------------------------------------------------------------------------------------
\begin{verbnobox}[\mbox{}\small]
 y                   | combine ............................ !
\end{verbnobox}
%----------------------------------------------------------------------------------------------------------------------------------
This parameter (\verb!combine!) enables the combination of the detection images of sources and filaments over wavebands; this
parameter may be changed to `\verb!n!', when \textsl{getsf} must be restarted to process only the other steps.
%----------------------------------------------------------------------------------------------------------------------------------
\begin{verbnobox}[\mbox{}\small]
 y 2                 | detect skelsig ..................... !
\end{verbnobox}
%----------------------------------------------------------------------------------------------------------------------------------
The 1st parameter (\verb!detect!) enables the detection of sources and filaments; this parameter may be changed to `\verb!n!', when
\textsl{getsf} must be restarted to process only the other steps. The 2nd parameter (\verb!skelsig!) is the desired significance of
the skeletons (crests of filaments); the default value may be changed, although this is not recommended.
%----------------------------------------------------------------------------------------------------------------------------------
\begin{verbnobox}[\mbox{}\small]
 y 1 1 50 n          | measure goodmin itn itx sproima .... !
\end{verbnobox}
%----------------------------------------------------------------------------------------------------------------------------------
The 1st parameter (\verb!measure!) enables the measurements of sources and filaments; this parameter may be changed to `\verb!n!',
when \textsl{getsf} must be restarted to process only the other steps. The 2nd parameter (\verb!goodmin!) defines the minimum
goodness of sources that are cataloged. When the measurements are enabled, but they must be restarted for any reason, the 3rd
parameter (\verb!itn!) defines the number of the iteration to restart with for the component of sources. The 4th parameter
(\verb!itx!) defines the maximum number of the measurement iterations (\verb!50!). The last parameter (\verb!sproima!), defaulting
to `\verb!n!', may be changed to `\verb!y!' to enable creation of small individual images and intensity profiles of each
background-subtracted and deblended cataloged source, as well as images of the sources in the original images.
%----------------------------------------------------------------------------------------------------------------------------------
\begin{verbnobox}[\mbox{}\small]
 y                   | visualize .......................... !
\end{verbnobox}
%----------------------------------------------------------------------------------------------------------------------------------
This parameter (\verb!visualize!) enables the visualization of extracted sources and filaments; this parameter may be changed to
`\verb!n!', when \textsl{getsf} must be restarted to process only the other steps. 
%The 2nd parameter (\verb!regcolor!) is the
%desired color of sources and skeletons in the \textsl{ds9} region files \verb!*.reg! created by \textsl{getsf} to aid users in
%visualizations and analyses. The \textsl{getsf} colors, abbreviated as \verb!red!, \verb!gre!, \verb!yel!, \verb!blu!, \verb!ora!, 
%\verb!pur!, \verb!cya!, \verb!mag!, \verb!lim!, \verb!pin!, \verb!tea!, \verb!lav!, \verb!bro!, \verb!bei!, \verb!mar!, 
%\verb!min!, 
%\verb!oli!, \verb!apr!, \verb!gra!, \verb!whi!, and \verb!bla!, somewhat differ from the standard colors to make them distinct.
%%\verb!red!, \verb!gre!, \verb!blu!, \verb!cya!, \verb!mag!, \verb!yel!, 
%%\verb!ora!, \verb!pin!, \verb!pur!, \verb!tea!, \verb!tan!, \verb!bro!, \verb!whi!, \verb!gra!, \verb!bla!.
%----------------------------------------------------------------------------------------------------------------------------------
\begin{verbnobox}[\mbox{}\small]
 5                   | nsigmacombos (!expert) ............. !
\end{verbnobox}
%----------------------------------------------------------------------------------------------------------------------------------
This parameter (\verb!nsigmacombos!) defines the default number of standard deviations in thresholding the component of sources,
when combining them over wavelengths. The default value (\verb!5!) is optimal and it should never be changed by non-experts.
%----------------------------------------------------------------------------------------------------------------------------------
\begin{verbnobox}[\mbox{}\small]
 2                   | nsigmacombof (!expert) ............. !
\end{verbnobox}
%----------------------------------------------------------------------------------------------------------------------------------
This parameter (\verb!nsigmacombof!) defines the default number of standard deviations in thresholding the component of filaments,
when combining them over wavelengths. The default value (\verb!2!) is optimal and it should never be changed by non-experts.
%----------------------------------------------------------------------------------------------------------------------------------
\begin{verbnobox}[\mbox{}\small]
 1                   | nwavesdetect ....................... !
\end{verbnobox}
%----------------------------------------------------------------------------------------------------------------------------------
This parameter (\verb!nwavesdetect!) defines the required minimum number of images (wavebands), in which a source is detected and
significant, to be included in the final catalog of sources. This value may be increased to 2, in order to reduce the number of
potentially spurious sources in the catalogs. However, higher values of this parameter may also lead to the elimination of the real
sources.
%----------------------------------------------------------------------------------------------------------------------------------
\begin{verbnobox}[\mbox{}\small]
 2 ?                 | savespace delete ................... !
\end{verbnobox}
%----------------------------------------------------------------------------------------------------------------------------------
The 1st parameter (\verb!savespace!) controls removal of the files created by \textsl{getsf} (when it has finished), defining the
hard drive cleaning level: \verb!0! (none), \verb!1! (medium), or \verb!2! (maximum). In the 2nd parameter (\verb!delete!), the
user can specify, whether \textsl{getsf} must wait for a confirmation before deleting the files (`\verb!?!') or assume that the
user would confirm (`\verb!y!') or disapprove (`\verb!n!') the actual deletion of files, without asking any question.
%----------------------------------------------------------------------------------------------------------------------------------
\begin{verbnobox}[\mbox{}\small]
 0                   | verbosity .......................... !
\end{verbnobox}
%----------------------------------------------------------------------------------------------------------------------------------
This parameter (\verb!verbosity!) defines three levels for the output produced by \textsl{getsf} on the user screen: `\verb!0!'
(minimum), `\verb!1!' (medium), or `\verb!2!' (maximum).

\begin{quote} %--------------------------------------------------------------------------------------------------------------------
\textbf{NOTE}. The maximum size (width) of sources (filaments) of interest (see above) can be conveniently estimated by opening
an image in \textsl{ds9} and creating a circular region, fully surrounding the \emph{entire} emission peak of the largest source
(filament) of interest; the \emph{radius} of the circular region (in arcsec) is the value that must replace the respective
placeholder.
\end{quote} %----------------------------------------------------------------------------------------------------------------------

%||||||||||||||||||||||||||||||||||||||||||||||||||||||||||||||||||||||||||||||||||||||||||||||||||||||||||||||||||||||||||||||||||
%||||||||||||||||||||||||||||||||||||||||||||||||||||||||||||||||||||||||||||||||||||||||||||||||||||||||||||||||||||||||||||||||||

\section{Extracting sources and filaments}
\label{extract}

This section presents examples of the \textsl{getsf} extraction runs, illustrating how \textsl{getsf} runs in sequential and
parallel modes. It is assumed that the prepared data set consists of the five \emph{Herschel} images, observed in the PACS/SPIRE
fast parallel mode, and an additional high-resolution surface density image. It is also assumed that the user has supplied all
required information about prepared images and the maximum sizes of interest in the upper section of \verb!+getsf.cfg!, for example:
%----------------------------------------------------------------------------------------------------------------------------------
\begin{verbnobox}[\mbox{}\small]
#_______________________________________________________________________________________________
#                                                                                               
# GETSF • Multi-Scale Multi-Wavelength Source & Filament Extraction • Version 241122 • Alexander
#_______________________________________________________________________________________________
#                    |                                      !
# VALUES             | PARAMETERS                           ! COMMENT  NB: default beams are for
#____________________|______________________________________!___________________________________
#                    |                                      !
 universe 1000       | prefix distance .................... ! prefix for naming files (never use
 6                   | nwmax ...............{ nwmax < 100 } ! maximum number of wavelengths spec
 6 1 2 3 4 5 6       | nwaves nw[i] ....................... ! number of waveblengths and their i
 070  8.4 17 25 n col | ................................... ! wave, beam, [-]srcmaxsize, [-]film
 ../../images ../070 | .................................... ! paths to image and monochromatic d
 universe_070        | .................................... ! prepared image name
 160 13.5 27 40 n col | ................................... ! wave, beam, [-]srcmaxsize, [-]film
 ../../images ../160 | .................................... ! paths to image and monochromatic d
 universe_160        | .................................... ! prepared image name
 165 13.5 27 40 n col | ................................... ! wave, beam, [-]srcmaxsize, [-]film
 ../../images/hires/+results ../165 | ..................... ! paths to image and monochromatic d
 hi.surface.density.r13p5 | ............................... ! prepared image name
 250 18.2 36 55 n col | ................................... ! wave, beam, [-]srcmaxsize, [-]film
 ../../images ../250 | .................................... ! paths to image and monochromatic d
 universe_250        | .................................... ! prepared image name
 350 24.9 50 75 n col | ................................... ! wave, beam, [-]srcmaxsize, [-]film
 ../../images ../350 | .................................... ! paths to image and monochromatic d
 universe_350        | .................................... ! prepared image name
 500 36.3 72 109 n col | .................................. ! wave, beam, [-]srcmaxsize, [-]film
 ../../images ../500 | .................................... ! paths to image and monochromatic d
 universe_500        | .................................... ! prepared image name
\end{verbnobox}
%----------------------------------------------------------------------------------------------------------------------------------

%||||||||||||||||||||||||||||||||||||||||||||||||||||||||||||||||||||||||||||||||||||||||||||||||||||||||||||||||||||||||||||||||||

\subsection{Running \textsl{getsf} sequentially for all images}
\label{sequence}

The bottom section of the configuration file tells \textsl{getsf} how to proceed with the extraction:
%----------------------------------------------------------------------------------------------------------------------------------
\begin{verbnobox}[\mbox{}\small]
 y y                 | sources filaments .................. ! extract sources and/or filaments {
 y 1 1 50            | separate its1 itf1 niter ........... ! separate components and background
 y 1 1               | flatten itsd1 itfd1 ................ ! flatten background-subtracted comp
 y                   | complete ........................... ! complete separation and flattening
 y                   | combine ............................ ! combine decomposed structures over
 y 2                 | detect skelsig ..................... ! detect structures; skeleton signif
 y 1 1 50 n          | measure goodmin itn itx sproima .... ! measure structures; minimum goodne
 y                   | visualize .......................... ! visualize sources
 5                   | nsigmacombos (!expert) ............. ! number of sigmas for combining sou
 2                   | nsigmacombof (!expert) ............. ! number of sigmas for combining fil
 1                   | nwavesdetect ....................... ! required minimum number of wavelen
 2 ?                 | savespace delete ................... ! space saving level {0|1|2}: clean 
 0                   | verbosity .......................... ! verbosity level {0|1|2} for the sc
#_______________________________________________________________________________________________
\end{verbnobox}
%----------------------------------------------------------------------------------------------------------------------------------
The first line tells \textsl{getsf} that both sources and filaments must be extracted. 
\begin{quote} %--------------------------------------------------------------------------------------------------------------------
\textbf{NOTE}. The correct approach to \textsl{getsf} extractions is first to inspect the observed images for presence of
filaments. If no filament-like structures are visible or there is no interest in the filaments, then processing of filaments must
be disabled by setting \verb!filaments! to `\verb!n!'. Processing of sources must \emph{always} be enabled: the source-like peaks
must be separated from filaments or background fluctuations. If the user is interested exclusively in the filament extraction, it
is possible to reduce the \textsl{getsf} execution time by limiting the source processing only to the source separation step. When
\textsl{getsf} has to be restarted from an intermediate step, sources may be disabled, depending on the circumstances.
\end{quote} %----------------------------------------------------------------------------------------------------------------------
If there is interest only in extracting sources, the \textsl{getsf} execution time can be reduced by roughly a factor of two, by
setting \verb!filaments! to `\verb!n!':
%----------------------------------------------------------------------------------------------------------------------------------
\begin{verbnobox}[\mbox{}\small]
 y n                 | sources filaments .................. ! extract sources and/or filaments {
\end{verbnobox}
%----------------------------------------------------------------------------------------------------------------------------------
The next block of parameters controls execution of the corresponding processing steps of \textsl{getsf}:
%----------------------------------------------------------------------------------------------------------------------------------
\begin{verbnobox}[\mbox{}\small]
 y 1 1 50            | separate its1 itf1 niter ........... ! separate components and background
 y 1 1               | flatten itsd1 itfd1 ................ ! flatten background-subtracted comp
 y                   | complete ........................... ! complete separation and flattening
 y                   | combine ............................ ! combine decomposed structures over
 y 2                 | detect skelsig ..................... ! detect structures; skeleton signif
 y 1 1 50 n          | measure goodmin itn itx sproima .... ! measure structures; minimum goodne
 y                   | visualize .......................... ! visualize sources
\end{verbnobox}
%----------------------------------------------------------------------------------------------------------------------------------
By default, they instruct \textsl{getsf} to process all images in the sequence of their appearance in the upper section of the
configuration file. This is the simplest approach to running \textsl{getsf}: the user makes the extraction directory (\verb!runs!)
the current one and starts the code within a new \textsl{screen} session:
%----------------------------------------------------------------------------------------------------------------------------------
\begin{verbnobox}[\mbox{}\small]
> cd 241122/runs
> screen -h 1000 -LS get
> getsf
\end{verbnobox}
%----------------------------------------------------------------------------------------------------------------------------------
Having started the code, the user may forget about it for some time, until it finishes (the time may be very long for many large
images). This simplest approach to running \textsl{getsf} is quite inefficient for the multi-wavelength extractions, because most
of the time the code spends in the first two steps (\verb!separate! and \verb!flatten!), processing the images completely
independently of each other. A natural and easy way to accelerate the extraction is to parallelize the image processing in
independent jobs, one \textsl{getsf} run per wavelength.

%||||||||||||||||||||||||||||||||||||||||||||||||||||||||||||||||||||||||||||||||||||||||||||||||||||||||||||||||||||||||||||||||||

\subsection{Running \textsl{getsf} in parallel for each image}
\label{parallel}

The parallel processing mode of images makes sense only for the steps of the separation of structural components (\verb!separate!),
flattening of background-subtracted images (\verb!flatten!), and completion of the two steps (\verb!complete!). This means that the
parameters \verb!combine!, \verb!detect!, \verb!measure!, and \verb!visualize! must be changed from `\verb!y!' to `\verb!n!', for
example:
%----------------------------------------------------------------------------------------------------------------------------------
\begin{verbnobox}[\mbox{}\small]
 y y                 | sources filaments .................. ! extract sources and/or filaments {
 y 1 1 50            | separate its1 itf1 niter ........... ! separate components and background
 y 1 1               | flatten itsd1 itfd1 ................ ! flatten background-subtracted comp
 y                   | complete ........................... ! complete separation and flattening
 n                   | combine ............................ ! combine decomposed structures over
 n 2                 | detect skelsig ..................... ! detect structures; skeleton signif
 n 1 1 50 n          | measure goodmin itn itx sproima .... ! measure structures; minimum goodne
 n                   | visualize .......................... ! visualize sources
\end{verbnobox}
%----------------------------------------------------------------------------------------------------------------------------------
Before running \textsl{getsf}, it is necessary to modify the line that defines the wavelengths to proceed. Edit the configuration
file by specifying the index of the 1st image:
%----------------------------------------------------------------------------------------------------------------------------------
\begin{verbnobox}[\mbox{}\small]
 1 1                 | nwaves nw[i] ....................... ! number of waveblengths and their i
\end{verbnobox}
%----------------------------------------------------------------------------------------------------------------------------------
Open a new \textsl{screen} session and start \textsl{getsf} for the 1st image:
%----------------------------------------------------------------------------------------------------------------------------------
\begin{verbnobox}[\mbox{}\small]
> screen -h 1000 -LS get1
> getsf
\end{verbnobox}
%----------------------------------------------------------------------------------------------------------------------------------
Edit the configuration file by specifying the index of the 2nd image:
%----------------------------------------------------------------------------------------------------------------------------------
\begin{verbnobox}[\mbox{}\small]
 1 2                 | nwaves nw[i] ....................... ! number of waveblengths and their i
\end{verbnobox}
%----------------------------------------------------------------------------------------------------------------------------------
Open a new \textsl{screen} session and start \textsl{getsf} for the 2nd image:
%----------------------------------------------------------------------------------------------------------------------------------
\begin{verbnobox}[\mbox{}\small]
> screen -h 1000 -LS get2
> getsf
\end{verbnobox}
%----------------------------------------------------------------------------------------------------------------------------------
Edit the configuration file by specifying the index of the 3rd image:
%----------------------------------------------------------------------------------------------------------------------------------
\begin{verbnobox}[\mbox{}\small]
 1 3                 | nwaves nw[i] ....................... ! number of waveblengths and their i
\end{verbnobox}
%----------------------------------------------------------------------------------------------------------------------------------
Open a new \textsl{screen} session and start \textsl{getsf} for the 3rd image:
%----------------------------------------------------------------------------------------------------------------------------------
\begin{verbnobox}[\mbox{}\small]
> screen -h 1000 -LS get3
> getsf
\end{verbnobox}
%----------------------------------------------------------------------------------------------------------------------------------
Edit the configuration file by specifying the index of the 4th image:
%----------------------------------------------------------------------------------------------------------------------------------
\begin{verbnobox}[\mbox{}\small]
 1 4                 | nwaves nw[i] ....................... ! number of waveblengths and their i
\end{verbnobox}
%----------------------------------------------------------------------------------------------------------------------------------
Open a new \textsl{screen} session and start \textsl{getsf} for the 4th image:
%----------------------------------------------------------------------------------------------------------------------------------
\begin{verbnobox}[\mbox{}\small]
> screen -h 1000 -LS get4
> getsf
\end{verbnobox}
%----------------------------------------------------------------------------------------------------------------------------------
Edit the configuration file by specifying the index of the 5th image:
%----------------------------------------------------------------------------------------------------------------------------------
\begin{verbnobox}[\mbox{}\small]
 1 5                 | nwaves nw[i] ....................... ! number of waveblengths and their i
\end{verbnobox}
%----------------------------------------------------------------------------------------------------------------------------------
Open a new \textsl{screen} session and start \textsl{getsf} for the 5th image:
%----------------------------------------------------------------------------------------------------------------------------------
\begin{verbnobox}[\mbox{}\small]
> screen -h 1000 -LS get5
> getsf
\end{verbnobox}
%----------------------------------------------------------------------------------------------------------------------------------
Edit the configuration file by specifying the index of the 6th image:
%----------------------------------------------------------------------------------------------------------------------------------
\begin{verbnobox}[\mbox{}\small]
 1 6                 | nwaves nw[i] ....................... ! number of waveblengths and their i
\end{verbnobox}
%----------------------------------------------------------------------------------------------------------------------------------
Open a new \textsl{screen} session and start \textsl{getsf} for the 6th image:
%----------------------------------------------------------------------------------------------------------------------------------
\begin{verbnobox}[\mbox{}\small]
> screen -h 1000 -LS get6
> getsf
\end{verbnobox}
%----------------------------------------------------------------------------------------------------------------------------------
Having started the parallel jobs, the user may forget about them for some time, until they all finish (the time may be very long
for very large images). Then it is necessary to invert the controls in the configuration file
(`\verb!n!'\,$\leftrightarrow$\,`\verb!y!') to tell \textsl{getsf} to complete all remaining steps:
%----------------------------------------------------------------------------------------------------------------------------------
\begin{verbnobox}[\mbox{}\small]
 y y                 | sources filaments .................. ! extract sources and/or filaments {
 n 1 1 50            | separate its1 itf1 niter ........... ! separate components and background
 n 1 1               | flatten itsd1 itfd1 ................ ! flatten background-subtracted comp
 n                   | complete ........................... ! complete separation and flattening
 y                   | combine ............................ ! combine decomposed structures over
 y 2                 | detect skelsig ..................... ! detect structures; skeleton signif
 y 1 1 50 n          | measure goodmin itn itx sproima .... ! measure structures; minimum goodne
 y                   | visualize .......................... ! visualize sources
\end{verbnobox}
%----------------------------------------------------------------------------------------------------------------------------------
The final \textsl{getsf} run completes the extraction (in a much shorter time than that required by the parallel steps):
%----------------------------------------------------------------------------------------------------------------------------------
\begin{verbnobox}[\mbox{}\small]
> screen -h 1000 -LS get
> getsf
\end{verbnobox}
%----------------------------------------------------------------------------------------------------------------------------------
 
%||||||||||||||||||||||||||||||||||||||||||||||||||||||||||||||||||||||||||||||||||||||||||||||||||||||||||||||||||||||||||||||||||

\subsection{Measuring filaments on all spatial scales}
\label{measureall}

Standard extractions with \textsl{getsf} provide measurements for only the filaments detected on the scales around the maximum
width of filaments of interest (set by the user in \verb!+getsf.cfg!), despite the fact that the method obtains skeletons for
filaments detected on every scale, down to the smallest one (around the angular resolution). To produce the filament measurements
on a smaller scale, it is necessary to determine and subtract the filament background appropriate for that scale. This can be
accomplished by an additional \textsl{getsf} run, in which the maximum width of \emph{filaments} in \verb!+getsf.cfg! is set to the
value of a smaller scale, for which skeletons were derived in the standard run (scale \verb!KKKK!, see Sect.~\ref{skeletons}). Only
the maximum size of filaments is to be changed, the one for sources must remain the same as in the standard extraction; using the
example from Sect.~\ref{extract}, only the \verb!YY! value must be changed:
%----------------------------------------------------------------------------------------------------------------------------------
\begin{verbnobox}[\mbox{}\small]
 165 13.5 27 YY n col | ................................... ! wave, beam, [-]srcmaxsize, [-]film
\end{verbnobox}
%----------------------------------------------------------------------------------------------------------------------------------
It is strongly advised to change control parameters in \verb!+getsf.cfg! to tell \textsl{getsf} that only a few precessing steps
are necessary:
%----------------------------------------------------------------------------------------------------------------------------------
\begin{verbnobox}[\mbox{}\small]
 n y                 | sources filaments .................. ! extract sources and/or filaments {
 y 1 1 50            | separate its1 itf1 niter ........... ! separate components and background
 n 1 1               | flatten itsd1 itfd1 ................ ! flatten background-subtracted comp
 n                   | complete ........................... ! complete separation and flattening
 n                   | combine ............................ ! combine decomposed structures over
 n 2                 | detect skelsig ..................... ! detect structures; skeleton signif
 y 1 1 50 n          | measure goodmin itn itx sproima .... ! measure structures; minimum goodne
 y                   | visualize .......................... ! visualize sources
\end{verbnobox}
%----------------------------------------------------------------------------------------------------------------------------------
The additional execution times are very short in comparison to the standard \textsl{getsf} run, because just a single
time-consuming processing step (filament separation) is completed. Such additional runs can be repeated for all scales found in the
skeletons subdirectory (Sect.~\ref{skeletons}).

%||||||||||||||||||||||||||||||||||||||||||||||||||||||||||||||||||||||||||||||||||||||||||||||||||||||||||||||||||||||||||||||||||
%||||||||||||||||||||||||||||||||||||||||||||||||||||||||||||||||||||||||||||||||||||||||||||||||||||||||||||||||||||||||||||||||||

\section{Understanding extraction output files}
\label{output}

Names of some output files contain numbers: \verb!WWW! is the waveband ($\mu$m), \verb!NN! the waveband number, \verb!JJ! the scale
number, \verb!OO! the angular resolution ($\arcsec$), \verb!TT! the iteration number, \verb!SSSSSS! the decomposition scale
($\arcsec$), \verb!FFFF! the maximum scale ($\arcsec$) of the partially reconstructed filament, \verb!GG! the skeleton
significance, \verb!XX! or \verb!XXX! the maximum size ($\arcsec$) of sources, \verb!YY! or \verb!YYY! the maximum size ($\arcsec$)
of filaments, \verb!KKKK! the skeleton scale ($\arcsec$), \verb!NNNNN! the cataloged source number, and \verb!NNNN! the filament
number. The images with \verb!obs! and \verb!det! in their names are related to the original images and detection images,
respectively.

%||||||||||||||||||||||||||||||||||||||||||||||||||||||||||||||||||||||||||||||||||||||||||||||||||||||||||||||||||||||||||||||||||

\subsection{Directory structure created by \textsl{getsf} runs}
\label{dirstructure}

When \textsl{getsf} is run in the extraction directory (\verb!241122/runs!) with the configuration example described in
Sect.~\ref{extract}, the following structure of sub-directories is created:\\
%----------------------------------------------------------------------------------------------------------------------------------
\vspace{-2mm}\noindent\\
{\small\verb!241122/070!} $\,\,\leadsto$ monochromatic processing (70 $\mu$m image)\\
{\small\verb!241122/160!} $\,\,\leadsto$ monochromatic processing (160 $\mu$m image)\\
{\small\verb!241122/165!} $\,\,\leadsto$ monochromatic processing (surface density image with the 160 $\mu$m image resolution)\\
{\small\verb!241122/250!} $\,\,\leadsto$ monochromatic processing (250 $\mu$m image)\\
{\small\verb!241122/350!} $\,\,\leadsto$ monochromatic processing (350 $\mu$m image)\\
{\small\verb!241122/500!} $\,\,\leadsto$ monochromatic processing (500 $\mu$m image)\\
{\small\verb!241122/components!} $\,\,\leadsto$ separated structural components (all wavebands)\\
{\small\verb!241122/components/+obsimage_partial!} $\,\,\leadsto$ partially reconstructed images of original images\\
{\small\verb!241122/components/+filaments_partial!} $\,\,\leadsto$ partially reconstructed images of filaments\\
{\small\verb!241122/runs/results!} $\,\,\leadsto$ detection and measurement results\\
{\small\verb!241122/runs/results/+components!} $\,\,\leadsto$ improved structural components (all wavebands)\\
{\small\verb!241122/runs/results/+sources/+catalogs!} $\,\,\leadsto$ catalogs of extracted sources\\
{\small\verb!241122/runs/results/+sources/+images!} $\,\,\leadsto$ individual images of extracted sources\\
{\small\verb!241122/runs/results/+sources/+iterations!} $\,\,\leadsto$ measurement iterations of sources (technical)\\
{\small\verb!241122/runs/results/+sources/+profiles!} $\,\,\leadsto$ individual profiles of extracted sources\\
{\small\verb!241122/runs/results/+sources/+visuals!} $\,\,\leadsto$ visualization of extracted sources\\
{\small\verb!241122/runs/results/+skeletons!} $\,\,\leadsto$ derived skeletons of filaments\\
{\small\verb!241122/runs/results/+filaments/+images/070!} $\,\,\leadsto$ measurements of filaments (images)\\
{\small\verb!241122/runs/results/+filaments/+images/160!} $\,\,\leadsto$ measurements of filaments (images)\\
{\small\verb!241122/runs/results/+filaments/+images/165!} $\,\,\leadsto$ measurements of filaments (images)\\
{\small\verb!241122/runs/results/+filaments/+images/250!} $\,\,\leadsto$ measurements of filaments (images)\\
{\small\verb!241122/runs/results/+filaments/+images/350!} $\,\,\leadsto$ measurements of filaments (images)\\
{\small\verb!241122/runs/results/+filaments/+images/500!} $\,\,\leadsto$ measurements of filaments (images)\\
{\small\verb!241122/runs/results/+filaments/+profiles/070!} $\,\,\leadsto$ measurements of filaments (catalogs)\\
{\small\verb!241122/runs/results/+filaments/+profiles/160!} $\,\,\leadsto$ measurements of filaments (catalogs)\\
{\small\verb!241122/runs/results/+filaments/+profiles/165!} $\,\,\leadsto$ measurements of filaments (catalogs)\\
{\small\verb!241122/runs/results/+filaments/+profiles/250!} $\,\,\leadsto$ measurements of filaments (catalogs)\\
{\small\verb!241122/runs/results/+filaments/+profiles/350!} $\,\,\leadsto$ measurements of filaments (catalogs)\\
{\small\verb!241122/runs/results/+filaments/+profiles/500!} $\,\,\leadsto$ measurements of filaments (catalogs)\\
{\small\verb!241122/runs/results/+filaments/+visuals!} $\,\,\leadsto$ visualization of extracted filaments
%----------------------------------------------------------------------------------------------------------------------------------
\begin{quote} %--------------------------------------------------------------------------------------------------------------------
\textbf{NOTE}. When an extraction has finished and the user has no plans to restart the same extraction in the future, it is
possible to save considerable disk space by removing unnecessary files. In the case that the extraction is indeed deemed final, the
user needs just the \texttt{components} and \verb!results! directories in the future. Therefore, it may be recommended that the
user removes all other directories: \verb!070!, \verb!160!, \verb!165!, \verb!250!, \verb!350!, \verb!500!, \verb!+iterations!, and
all \emph{files} from the extraction directory \verb!runs!, \emph{with an exception} of the \textsl{getsf} configuration and log
files (\verb!+getsf.cfg!, \verb!+log.*!) that must be preserved as a documentation. Another possibility is to compress the
directories using an archiver that produces highly compressed solid archives, e.g., \textsl{rar}
(\url{https://www.rarlab.com/download.htm}).
\end{quote} %----------------------------------------------------------------------------------------------------------------------

%||||||||||||||||||||||||||||||||||||||||||||||||||||||||||||||||||||||||||||||||||||||||||||||||||||||||||||||||||||||||||||||||||

\subsection{Directory for \textsl{getsf} extraction runs}
\label{extractdir}

The directory \verb!runs! contains the following files:\\

\vspace{-2mm}\noindent
{\small\verb!<fieldname>.WWW.obs.sources.n.fits!} $\,\,\leadsto$ regularized component of sources\\
{\small\verb!<fieldname>.WWW.obs.sources.n.cb.fits!} $\,\,\leadsto$ interpolated background of regularized component of sources\\
{\small\verb!<fieldname>.WWW.obs.filaments.n.fits!} $\,\,\leadsto$ regularized component of filaments\\
{\small\verb!<fieldname>.WWW.det.sflatfactor.fits!} $\,\,\leadsto$ flattening image for the component of sources\\
{\small\verb!<fieldname>.WWW.det.fflatfactor.fits!} $\,\,\leadsto$ flattening image for the component of filaments\\
{\small\verb!<fieldname>.WWW.det.fits!} $\,\,\leadsto$ detection image for the component of sources or filaments\\
{\small\verb!<fieldname>.WWW.obs.fits!} $\,\,\leadsto$ original observed image\\
{\small\verb!<fieldname>.WWW.omask.fits!} $\,\,\leadsto$ observational mask\\
{\small\verb!<fieldname>.WWW.bs.omask.fits!} $\,\,\leadsto$ observational mask for background subtraction\\
{\small\verb!<fieldname>.f.det.JJ.cSSSSSSas.an.fits!} $\,\,\leadsto$ detection images for filaments combined over waves\\
{\small\verb!<fieldname>.s.det.JJ.cSSSSSSas.an.fits!} $\,\,\leadsto$ detection images for sources combined over waves\\
{\small\verb!<fieldname>.s.det.JJ.cSSSSSSas.as.fits!} $\,\,\leadsto$ detection images for sources combined over waves\\
{\small\verb!<fieldname>.f.coadded.omask.fits!} $\,\,\leadsto$ observational mask accumulated over wavebands for filaments\\
{\small\verb!<fieldname>.s.coadded.omask.fits!} $\,\,\leadsto$ observational mask accumulated over wavebands for sources\\
{\small\verb!<fieldname>.combining.WWW.omask.fits!} $\,\,\leadsto$ special mask (optional) to exclude bad pixels in combining\\
{\small\verb!sm.<fieldname>.WWW.obs.subsfp.fits!} $\,\,\leadsto$ footprints of sub-structured sources (if any)\\
{\small\verb!sm.<fieldname>.s.det.footprints.fits!} $\,\,\leadsto$ footprints of all detected sources\\
{\small\verb!sm.<fieldname>.s.det.unrotated.fits!} $\,\,\leadsto$ positions of all detected sources\\
{\small\verb!sm.<fieldname>.s.det.wcs!} $\,\,\leadsto$ conversion table from pixels to world coordinates of sources\\
{\small\verb!sm.<fieldname>.s.det.cat!} $\,\,\leadsto$ detection catalog of sources\\
{\small\verb!sm.<fieldname>.WWW.obs.smf.cat!} $\,\,\leadsto$ catalog of the detected sources updated in measurement iterations\\
{\small\verb!sm.seg.<fieldname>.s.det.JJ.cSSSSSSas.fits!} $\,\,\leadsto$ segmentation of sources at scale \verb!SSSSSS!\\
{\small\verb!sm.seg.<fieldname>.s.det.JJ.cSSSSSSas.fp.fits!} $\,\,\leadsto$ footprints of detected sources at scale \verb!SSSSSS!\\
{\small\verb!sm.seg.<fieldname>.s.det.JJ.cSSSSSSas.cat!} $\,\,\leadsto$ detection catalog of sources at scale \verb!SSSSSS!\\
{\small\verb!bigfoot.WWW.fits!} $\,\,\leadsto$ annuli around the source footprints (used for flux errors estimation)\\
{\small\verb!flatfoot.WWW.fits!} $\,\,\leadsto$ converged footprints of all acceptable sources\\
{\small\verb!<fieldname>.mask.s!$\sim$\verb!max.wWWW.sum.hi.fits!} $\,\,\leadsto$ masks of the structures accumulated over spatial scales\\
{\small\verb!<fieldname>.mask.s!$\sim$\verb!max.wWWW!$\sim$\verb!WWW.sum.hi.fits!} $\,\,\leadsto$ masks accumulated over scales and wavebands\\
{\small\verb!<fieldname>.s.det.combine!} $\,\,\leadsto$ technical info on the detection images for sources being combined\\
{\small\verb!<fieldname>.f.det.combine!} $\,\,\leadsto$ technical info on the detection images for filaments being combined\\
{\small\verb!<fieldname>.s.det.combined!} $\,\,\leadsto$ technical info on the combined detection images for sources\\
{\small\verb!<fieldname>.f.det.combined!} $\,\,\leadsto$ technical info on the combined detection images for filaments\\
{\small\verb!<fieldname>.WWW.srcmaxsize!} $\,\,\leadsto$ technical info on the maximum size of sources\\
{\small\verb!<fieldname>.WWW.det.lastbright!} $\,\,\leadsto$ technical info on the last bright spatial scale\\
{\small\verb!+wave.WWW.done!} $\,\,\leadsto$ technical info on convergence of the source measurement iterations\\
{\small\verb!+wave.WWW.foot!} $\,\,\leadsto$ technical info on footprint expansion or shrinkage in measurement iterations

%||||||||||||||||||||||||||||||||||||||||||||||||||||||||||||||||||||||||||||||||||||||||||||||||||||||||||||||||||||||||||||||||||

\subsection{Directory for monochromatic processing}
\label{monochrom}

The directories \verb!runs/../WWW! contain the following files:\\

\vspace{-2mm}\noindent
{\small\verb!<fieldname>.WWW.obs.sources.fits!} $\,\,\leadsto$ separated component of sources\\
{\small\verb!<fieldname>.WWW.obs.filaments.fits!} $\,\,\leadsto$ separated component of filaments\\
{\small\verb!<fieldname>.WWW.obs.sources.n.fits!} $\,\,\leadsto$ regularized component of sources\\
{\small\verb!<fieldname>.WWW.obs.filaments.n.fits!} $\,\,\leadsto$ regularized component of filaments\\
{\small\verb!<fieldname>.WWW.det.sources.fits!} $\,\,\leadsto$ detection image of the separated component of sources\\
{\small\verb!<fieldname>.WWW.det.filaments.fits!} $\,\,\leadsto$ detection image of the separated component of filaments\\
{\small\verb!<fieldname>.WWW.obs.fbackground.fits!} $\,\,\leadsto$ background of the separated component of filaments\\
{\small\verb!<fieldname>.WWW.obs.sbackground.fits!} $\,\,\leadsto$ background of the separated component of sources\\
{\small\verb!<fieldname>.WWW.det.fits!} $\,\,\leadsto$ updated background image of the separated component\\
{\small\verb!<fieldname>.WWW.omask.fits!} $\,\,\leadsto$ observational mask image\\
{\small\verb!<fieldname>.WWW.det.decompose!} $\,\,\leadsto$ technical info on the decomposed single scale images\\
{\small\verb!<fieldname>.WWW.det.lastbright!} $\,\,\leadsto$ technical info on last bright spatial scale\\
{\small\verb!<fieldname>.WWW.det.separateso!} $\,\,\leadsto$ technical info on the decomposed images of sources\\
{\small\verb!<fieldname>.WWW.det.separatefi!} $\,\,\leadsto$ technical info on the decomposed images of filaments\\
{\small\verb!+log.GETSF.iter.NN!} $\,\,\leadsto$ technical info on the threshold iterations for decomposed images

\subsubsection{Directory for separation of sources or filaments}
\label{separate}

The directories \verb!runs/../WWW/+components/+{sources|filaments}! contain the following files:\\

\vspace{-2mm}\noindent
{\small\verb!<fieldname>.WWW.det.bg.TT.fits!} $\,\,\leadsto$ background of the component in iterations\\
{\small\verb!<fieldname>.WWW.det.co.TT.fits!} $\,\,\leadsto$ separated component in iterations\\
{\small\verb!<fieldname>.WWW.det.bg.TT.m.fits!} $\,\,\leadsto$ background after subtraction of a constant\\
{\small\verb!<fieldname>.WWW.det.bg.TT.r.fits!} $\,\,\leadsto$ relative correction to the background in iterations\\
{\small\verb!<fieldname>.WWW.det.bg.fits!} $\,\,\leadsto$ background of the component at the latest iteration\\
{\small\verb!<fieldname>.WWW.det.co.fits!} $\,\,\leadsto$ separated component at the latest iteration\\
{\small\verb!<fieldname>.WWW.det.fits!} $\,\,\leadsto$ initial image for the component separation\\
{\small\verb!<fieldname>.WWW.det.fm.01.fits!} $\,\,\leadsto$ cleaning masks accumulated over scales at the first iteration\\
{\small\verb!<fieldname>.WWW.det.b.fits!} $\,\,\leadsto$ approximately reconstructed background of the component\\
{\small\verb!<fieldname>.WWW.obs.bg.sXXas.fits!} $\,\,\leadsto$ final separated background of the component\\
{\small\verb!<fieldname>.WWW.obs.bg.sXXas.n.fits!} $\,\,\leadsto$ final regularized background of the component\\
{\small\verb!<fieldname>.WWW.obs.bgs.sXXas.fits!} $\,\,\leadsto$ final separated component\\
{\small\verb!<fieldname>.WWW.obs.bgs.sXXas.std.fits!} $\,\,\leadsto$ standard deviations in the component\\
{\small\verb!<fieldname>.WWW.obs.bgs.sXXas.flat.fits!} $\,\,\leadsto$ flattened separated component\\
{\small\verb!<fieldname>.WWW.obs.bgs.sXXas.flat.std.fits!} $\,\,\leadsto$ standard deviations in the flat component\\
{\small\verb!<fieldname>.WWW.obs.bgs.sXXas.n.fits!} $\,\,\leadsto$ regularized component\\
{\small\verb!<fieldname>.WWW.obs.bgs.sXXas.n.flat.fits!} $\,\,\leadsto$ regularized flattened component\\
{\small\verb!<fieldname>.WWW.obs.bgs.sXXas.n.flat.std.fits!} $\,\,\leadsto$ standard deviations in the regularized component\\
{\small\verb!<fieldname>.WWW.obs.fits!} $\,\,\leadsto$ original observed image\\
{\small\verb!<fieldname>.WWW.det.sigma!} $\,\,\leadsto$ technical info on the standard deviation over all scales\\
{\small\verb!<fieldname>.WWW.det.decompose!} $\,\,\leadsto$ technical info on the decomposed images\\
{\small\verb!<fieldname>.WWW.det.lastbright!} $\,\,\leadsto$ technical info on the last bright spatial scale\\
{\small\verb!<fieldname>.WWW.det.separatebg!} $\,\,\leadsto$ technical info on the decomposed component\\
{\small\verb!<fieldname>.WWW.filmaxsize!} $\,\,\leadsto$ technical info on the maximum size of filaments

\subsubsection{Directory for flattening of sources or filaments}
\label{flatten}

The directories \verb!runs/../WWW/+components/+{srcfactor|filfactor}! contain the following files:\\

\vspace{-2mm}\noindent
{\small\verb!<fieldname>.WWW.det.bg.TT.fits!} $\,\,\leadsto$ background of the standard deviations in iterations\\
{\small\verb!<fieldname>.WWW.det.co.TT.fits!} $\,\,\leadsto$ separated component in iterations\\
{\small\verb!<fieldname>.WWW.det.bg.TT.m.fits!} $\,\,\leadsto$ background after subtraction of a constant\\
{\small\verb!<fieldname>.WWW.det.bg.TT.r.fits!} $\,\,\leadsto$ relative correction to the background in iterations\\
{\small\verb!<fieldname>.WWW.det.bg.fits!} $\,\,\leadsto$ background of standard deviations at the latest iteration\\
{\small\verb!<fieldname>.WWW.det.co.fits!} $\,\,\leadsto$ separated components at the latest iteration\\
{\small\verb!<fieldname>.WWW.det.fits!} $\,\,\leadsto$ initial image for the component separation\\
{\small\verb!<fieldname>.WWW.det.fm.01.fits!} $\,\,\leadsto$ cleaning masks accumulated over scales at the first iteration\\
{\small\verb!<fieldname>.WWW.det.b.fits!} $\,\,\leadsto$ approximately reconstructed background of the standard deviations\\
{\small\verb!<fieldname>.WWW.obs.bgs.sXXas.std.bg.sOOas.fits!} $\,\,\leadsto$ separated background (flattening image)\\
{\small\verb!<fieldname>.WWW.obs.fits!} $\,\,\leadsto$ initial image of standard deviations\\
{\small\verb!<fieldname>.WWW.det.sigma!} $\,\,\leadsto$ technical info on the standard deviation over all scales\\
{\small\verb!<fieldname>.WWW.det.decompose!} $\,\,\leadsto$ technical info on the decomposed images\\
{\small\verb!<fieldname>.WWW.det.lastbright!} $\,\,\leadsto$ technical info on the last bright spatial scale\\
{\small\verb!<fieldname>.WWW.det.separateff!} $\,\,\leadsto$ technical info on the decomposed component\\
{\small\verb!<fieldname>.WWW.filmaxsize!} $\,\,\leadsto$ technical info on the maximum size of filaments

\subsubsection{Directory for completion of sources or filaments}
\label{complete}

These directories (\verb!runs/../WWW/+components/+{scomplete|fcomplete}!) contain the following files:\\

\vspace{-2mm}\noindent
{\small\verb!<fieldname>.WWW.det.fits!} $\,\,\leadsto$ flattened detection image for the structural component\\
{\small\verb!<fieldname>.WWW.det.sigma!} $\,\,\leadsto$ technical info on the standard deviation over all scales

%||||||||||||||||||||||||||||||||||||||||||||||||||||||||||||||||||||||||||||||||||||||||||||||||||||||||||||||||||||||||||||||||||

\subsection{Directory for separated structural components}
\label{components}

This directory \verb!runs/../components! contains the following files:\\

\vspace{-2mm}\noindent
{\small\verb!<fieldname>.WWW.det.sources.fits!} $\,\,\leadsto$ separated component of sources (flattened)\\
{\small\verb!<fieldname>.WWW.det.filaments.sXXas.fits!} $\,\,\leadsto$ separated component of filaments (flattened)\\
{\small\verb!<fieldname>.WWW.det.sflatfactor.fits!} $\,\,\leadsto$ flattening factor image for the component of sources\\
{\small\verb!<fieldname>.WWW.det.fflatfactor.sXXas.fits!} $\,\,\leadsto$ flattening factor image for the component of filaments\\
{\small\verb!<fieldname>.WWW.obs.sources.fits!} $\,\,\leadsto$ separated component of sources (observed)\\
{\small\verb!<fieldname>.WWW.obs.sources.n.fits!} $\,\,\leadsto$ regularized component of sources (observed)\\
{\small\verb!<fieldname>.WWW.obs.filaments.sXXas.fits!} $\,\,\leadsto$ separated component of filaments (observed)\\
{\small\verb!<fieldname>.WWW.obs.filaments.n.sXXas.fits!} $\,\,\leadsto$ regularized component of filaments (observed)\\
{\small\verb!<fieldname>.WWW.obs.sbackground+sources.fits!} $\,\,\leadsto$ background of sources plus sources\\
{\small\verb!<fieldname>.WWW.obs.fbackground.sXXas+filaments.sXXas.fits!} $\,\,\leadsto$ background of filaments plus filaments

\subsubsection{Directory for partially reconstructed images}
\label{partial}

The directories \verb!runs/../WWW/+components/+{obsimage|filaments}_partial! contain the following files:\\

\vspace{-2mm}\noindent
{\small\verb!<fieldname>.WWW.det.FFFFas.fits!} $\,\,\leadsto$ observed image or the filaments reconstructed up to scale \verb!FFFF!

%||||||||||||||||||||||||||||||||||||||||||||||||||||||||||||||||||||||||||||||||||||||||||||||||||||||||||||||||||||||||||||||||||

\subsection{Directory for output images and catalogs}
\label{results}

The directory \verb!runs/results! contains several subdirectories.

\subsubsection{Directory for final improved structural components}
\label{pluscompon}

The directory \verb!runs/results/+components! contains the following files:\\

\vspace{-2mm}\noindent
{\small\verb!<fieldname>.WWW.obs.fits!} $\,\,\leadsto$ original observed image\\
{\small\verb!<fieldname>.s.WWW.obs.sources.fits!} $\,\,\leadsto$ sources obtained by subtraction of their interpolated background\\
{\small\verb!<fieldname>.s.WWW.obs.sbackground.fits!} $\,\,\leadsto$ background of sources interpolated under their footprints\\
{\small\verb!<fieldname>.WWW.obs.fbackground.fits!} $\,\,\leadsto$ regularized background of separated component of filaments\\
{\small\verb!<fieldname>.WWW.obs.filaments.fits!} $\,\,\leadsto$ observed filaments recomputed as \verb!sbackground-fbackground!\\
{\small\verb!<fieldname>.WWW.det.filaments.fits!} $\,\,\leadsto$ flattened image of observed recomputed filaments

\subsubsection{Directory for catalogs of extracted sources}
\label{srccats}

The directory \verb!runs/results/+sources/+catalogs! contains the following files:\\

\vspace{-2mm}\noindent
{\small\verb!<fieldname>.s.WWW.sources.cat!} $\,\,\leadsto$ final catalog of extracted sources (numbering from detection)\\
{\small\verb!<fieldname>.s.WWW.sources.ok.cat!} $\,\,\leadsto$ final catalog of extracted sources (resorted and renumbered)\\
{\small\verb!<fieldname>.s.WWW.sources.ok.cat.WWW.good!} $\,\,\leadsto$ final catalog of only selected good sources\\
{\small\verb!<fieldname>.s.WWW.sources.ok.cat.WWW.ell.fk5.reg!} $\,\,\leadsto$ \textsl{ds9} region file of half-maximum ellipses\\
{\small\verb!<fieldname>.s.WWW.sources.ok.cat.WWW.ell.phy.reg!} $\,\,\leadsto$ \textsl{ds9} region file of half-maximum ellipses\\
{\small\verb!<fieldname>.s.WWW.sources.ok.cat.WWW.foo.fk5.reg!} $\,\,\leadsto$ \textsl{ds9} region file of footprint ellipses\\
{\small\verb!<fieldname>.s.WWW.sources.ok.cat.WWW.foo.phy.reg!} $\,\,\leadsto$ \textsl{ds9} region file of footprint ellipses\\
{\small\verb!<fieldname>.s.WWW.sources.ok.add.cat!} $\,\,\leadsto$ additional catalog for sources (resorted and renumbered)\\
{\small\verb!<fieldname>.s.WWW.obs.chisq.dat!} $\,\,\leadsto$ parameters of elliptical Gaussians fits to extracted sources\\
{\small\verb!sm.<fieldname>.s.det.cat!} $\,\,\leadsto$ catalog of detected sources

\subsubsection{Directory for images of individual sources}
\label{srccats}

The directory \verb!runs/results/+sources/+images! contains the following files:\\

\vspace{-2mm}\noindent
{\small\verb!sm.WWW.bgs.fs.dbl.src.okNNNNN.noMMMMM.fits!} $\,\,\leadsto$ images of background-subtracted and deblended sources\\
{\small\verb!sm.WWW.obs.src.okNNNNN.noMMMMM.fits!} $\,\,\leadsto$ observed images of extracted sources

\subsubsection{Directory for measurement iterations of sources}
\label{srciters}

The directory \verb!runs/results/+sources/+iterations! contains the following files:\\

\vspace{-2mm}\noindent
{\small\verb!<fieldname>.s.sources.TT.cat!} $\,\,\leadsto$ final catalog of extracted sources in measurement iterations\\
{\small\verb!<fieldname>.s.sources.ok.TT.cat!} $\,\,\leadsto$ final catalog of extracted sources (renumbered) in iterations\\
{\small\verb!<fieldname>.WWW.det.NN.cSSSSSSas.fpeak.deblending.dat!} $\,\,\leadsto$ info on single-scale source peak deblending\\
{\small\verb!sm.s.WWW.foot.expansion.TT.dat!} $\,\,\leadsto$ footprint expansion and shrinkage in measurement iterations\\
{\small\verb!sm.s.WWW.fpeak.deblending.TT.dat!} $\,\,\leadsto$ deblending source peak intensity in measurement iterations\\
{\small\verb!sm.<fieldname>.WWW.obs.foots.TT.fits!} $\,\,\leadsto$ footprints of sources in measurement iterations\\
{\small\verb!sm.s.<fieldname>.WWW.obs.smf.TT.cat!} $\,\,\leadsto$ catalog of detected sources updated in measurement iterations\\
{\small\verb!sm.<fieldname>.WWW.obs.smf.cat!} $\,\,\leadsto$ catalog of detected sources at the last measurement iteration\\
{\small\verb!sm.<fieldname>.WWW.obs.smf.previous.cat!} $\,\,\leadsto$ catalog of detected sources at previous to last iteration

\subsubsection{Directory for intensity profiles of sources}
\label{srccats}

The directory \verb!runs/results/+sources/+profiles! contains the following files:\\

\vspace{-2mm}\noindent
{\small\verb!sm.intens.pro.okNNNNN.noMMMMM.dat!} $\,\,\leadsto$ background-subtracted and deblended intensity profiles of sources

\subsubsection{Directory for visualization of extracted sources}
\label{srcvisuals}

The directory \verb!runs/results/+sources/+visuals! contains the following files:\\

\vspace{-2mm}\noindent
{\small\verb!<fieldname>.s.WWW.fitshapes.bground.fits!} $\,\,\leadsto$ background of sources by subtraction of Gaussian shapes\\
{\small\verb!<fieldname>.s.WWW.fitshapes.res.mean.fits!} $\,\,\leadsto$ mean difference of sources from Gaussian shapes\\
{\small\verb!<fieldname>.s.WWW.fitshapes.res.true.fits!} $\,\,\leadsto$ actual difference of sources from Gaussian shapes\\
{\small\verb!<fieldname>.s.WWW.fitshapes.sources.fits!} $\,\,\leadsto$ elliptical Gaussian shapes fitted to extracted sources\\
{\small\verb!<fieldname>.s.WWW.obs.cb.fits!} $\,\,\leadsto$ background of sources obtained by interpolation under footprints\\
{\small\verb!<fieldname>.s.WWW.obs.bs.fits!} $\,\,\leadsto$ sources obtained by subtraction of interpolated background\\
{\small\verb!<fieldname>.s.WWW.obs.cb.smo.fits!} $\,\,\leadsto$ interpolated background smoothed to the lowest resolution\\
{\small\verb!<fieldname>.s.WWW.obs.sources.n.cb.fits!} $\,\,\leadsto$ interpolated background in regularized component of sources\\
{\small\verb!<fieldname>.s.WWW.obs.fits!} $\,\,\leadsto$ original observed image\\
{\small\verb!<fieldname>.s.sources.ell.WWW.obs.fits!} $\,\,\leadsto$ half-maximum ellipses overlaid on observed image\\
{\small\verb!<fieldname>.s.sources.ell.WWW.obs.src.fits!} $\,\,\leadsto$ half-maximum ellipses on component of sources\\
{\small\verb!<fieldname>.s.sources.ell.WWW.obs.src.n.fits!} $\,\,\leadsto$ half-maximum ellipses on regularized sources\\
{\small\verb!<fieldname>.s.sources.ell.WWW.obs.fil.fits!} $\,\,\leadsto$ half-maximum ellipses on component of filaments\\
{\small\verb!<fieldname>.s.sources.ell.WWW.obs.fil+ske.global.fits!} $\,\,\leadsto$ half-max ellipses on global skeletons\\
{\small\verb!<fieldname>.s.sources.foo.WWW.obs.fits!} $\,\,\leadsto$ footprint ellipses overlaid on observed image\\
{\small\verb!<fieldname>.s.sources.foo.WWW.obs.src.fits!} $\,\,\leadsto$ footprint ellipses on component of sources\\
{\small\verb!<fieldname>.s.sources.foo.WWW.obs.src.n.fits!} $\,\,\leadsto$ footprint ellipses on regularized sources\\
{\small\verb!<fieldname>.s.sources.foo.WWW.obs.fil.fits!} $\,\,\leadsto$ footprint ellipses on component of filaments\\
{\small\verb!<fieldname>.s.sources.foo.all.fits!} $\,\,\leadsto$ footprint ellipses accumulated over all wavebands\\
{\small\verb!<fieldname>.s.WWW.det.foots.fits!} $\,\,\leadsto$ initial footprints of detected sources\\
{\small\verb!<fieldname>.s.WWW.obs.foots.fits!} $\,\,\leadsto$ final footprints of measured sources

\subsubsection{Directory for derived skeletons of filaments}
\label{skeletons}

The directory \verb!runs/results/+skeletons! contains the following files:\\

\vspace{-2mm}\noindent
{\small\verb!<fieldname>.f.det.ske.KKKKas.fits!} $\,\,\leadsto$ accumulated single-scale skeletons on scales around \verb!KKKK!\\
{\small\verb!<fieldname>.f.det.ske.KKKKas.sGG.fits!} $\,\,\leadsto$ skeletons with significance \verb!GG! obtained on scale \verb!KKKK!

\subsubsection{Directory for derived global skeletons of filaments}
\label{globals}

The directory \verb!runs/results/+skeletons/+globals! contains the following files:\\

\vspace{-2mm}\noindent
{\small\verb!<fieldname>.f.det.skeletons.KKKKas.sGG.fits!} $\,\,\leadsto$ isolated skeletons with significance \verb!GG! on scale \verb!KKKK!\\
{\small\verb!<fieldname>.f.det.skeletons.KKKKas.sGG.cat!} $\,\,\leadsto$ coordinates of the points of isolated skeletons\\
{\small\verb!<fieldname>.f.det.skeletons.KKKKas.sGG.reg!} $\,\,\leadsto$ \textsl{ds9} region file of isolated skeletons\\
{\small\verb!<fieldname>.f.det.skeletons.global.sGG.fits!} $\,\,\leadsto$ global skeletons on all spatial scales\\
{\small\verb!<fieldname>.f.det.skeletons.global.sGG.cat!} $\,\,\leadsto$ coordinates of the global skeletons\\
{\small\verb!<fieldname>.f.det.skeletons.global.sGG.reg!} $\,\,\leadsto$ \textsl{ds9} region file of global skeletons

\subsubsection{Directory for images used in skeleton derivation}
\label{sdetect}

The directory \verb!runs/results/+sdetect! contains the following files:\\

\vspace{-2mm}\noindent
{\small\verb!<fieldname>.f.det.fil.KKKKas.fits!} $\,\,\leadsto$ single-scale filament images accumulated on scales around \verb!KKKK!

\subsubsection{Directory for images from filament measurements}
\label{filimages}

The directory \verb!runs/results/+filaments/maxYYYas/+images/WWW! contains the following files:\\

\vspace{-2mm}\noindent
{\small\verb!<fieldname>.f.WWW.obs.ske.KKKKas.sGG!$\sim$\verb!angles.fits!} $\,\,\leadsto$ position angles at each skeleton point\\
{\small\verb!<fieldname>.f.WWW.obs.ske.KKKKas.sGG!$\sim$\verb!foots.fits!} $\,\,\leadsto$ two-sided footprints of filaments\\
{\small\verb!<fieldname>.f.WWW.obs.ske.KKKKas.sGG!$\sim$\verb!lindens.fits!} $\,\,\leadsto$ linear densities along crests of filaments\\
{\small\verb!<fieldname>.f.WWW.obs.ske.KKKKas.sGG!$\sim$\verb!normals.fits!} $\,\,\leadsto$ normals to skeletons at each point\\
{\small\verb!<fieldname>.f.WWW.obs.ske.KKKKas.sGG!$\sim$\verb!widths.fits!} $\,\,\leadsto$ average widths of filaments

\subsubsection{Directory for measurement catalogs of filaments}
\label{filprofiles}

The directory \verb!runs/results/+filaments/maxYYYas/+profiles/WWW! contains the following files:\\

\vspace{-2mm}\noindent
{\small\verb!<fieldname>.f.WWW.obs.ske.KKKKas.sGG!$\sim$\verb!coo.hires.cat!} $\,\,\leadsto$ coordinates of skeleton points\\
{\small\verb!<fieldname>.f.WWW.obs.ske.KKKKas.sGG!$\sim$\verb!coo.hires.reg!} $\,\,\leadsto$ \textsl{ds9} region file of skeletons\\
{\small\verb!<fieldname>.f.WWW.obs.ske.KKKKas.sGG!$\sim$\verb!lindens.cat!} $\,\,\leadsto$ average linear densities and masses\\
{\small\verb!<fieldname>.f.WWW.obs.ske.KKKKas.sGG!$\sim$\verb!lindens.fi.NNNN.cat!} $\,\,\leadsto$ linear densities along fil \verb!NNNN!\\
{\small\verb!<fieldname>.f.WWW.obs.ske.KKKKas.sGG!$\sim$\verb!profiles.fi.NNNN.mea.cat!} $\,\,\leadsto$ mean profiles for \verb!NNNN!\\
{\small\verb!<fieldname>.f.WWW.obs.ske.KKKKas.sGG!$\sim$\verb!profiles.fi.NNNN.med.cat!} $\,\,\leadsto$ median profiles for \verb!NNNN!\\
{\small\verb!<fieldname>.f.WWW.obs.ske.KKKKas.sGG!$\sim$\verb!profiles.fi.NNNN.sel.cat!} $\,\,\leadsto$ selected profiles for \verb!NNNN!\\
{\small\verb!<fieldname>.f.WWW.obs.ske.KKKKas.sGG!$\sim$\verb!selwidth.cat!} $\,\,\leadsto$ filament widths along selected normals\\
{\small\verb!<fieldname>.f.WWW.obs.ske.KKKKas.sGG!$\sim$\verb!widths.cat!} $\,\,\leadsto$ widths averaged along filaments

\subsubsection{Directory for visualization of skeletons on filaments}
\label{filvisuals}

The directory \verb!runs/results/+filaments/maxYYYas/+visuals! contains the following files:\\

\vspace{-2mm}\noindent
{\small\verb!<fieldname>.f.WWW.obs.filaments+ske.KKKKas.sGG.fits!} $\,\,\leadsto$ skeletons at \verb!KKKK! with \verb!GG! on filaments\\
{\small\verb!<fieldname>.f.WWW.obs.filaments+ske.KKKKas.sGG.reg!} $\,\,\leadsto$ \textsl{ds9} region file of skeletons\\
{\small\verb!<fieldname>.f.WWW.obs.filaments+ske.global.JJJJ!$\sim$\verb!KKKKas.fits!} $\,\,\leadsto$ global skeletons on filaments\\
{\small\verb!<fieldname>.f.WWW.obs.filaments+ske.global.JJJJ!$\sim$\verb!KKKKas.reg!} $\,\,\leadsto$ \textsl{ds9} region file of global skeletons

\subsubsection{Directory for visualization of filaments}
\label{filvisuals}

The directory \verb!runs/results/+filaments/maxXXas/+visuals/+nosmall! contains the following files:\\

\vspace{-2mm}\noindent
{\small\verb!<fieldname>.f.WWW.obs.filaments.scale.KKKKas.fits!} $\,\,\leadsto$ filaments with scales below \verb!KKKK! removed

\section{User agreement}
\label{agreement}

In what follows, \emph{this software} refers to the scripts \textsl{getsf}, \textsl{comfun}, \textsl{compile}, \textsl{convolve},
\textsl{hires}, \textsl{installg}, \textsl{iospeed}, \textsl{prepobs}, \textsl{resample}, and the associated utilities
\textsl{cleanbg}, \textsl{ellipse}, \textsl{expanda}, \textsl{extractx}, \textsl{fftconv}, \textsl{finalcat}, \textsl{fitfluxes},
\textsl{fixskels}, \textsl{fmeasure}, \textsl{imgstat}, \textsl{makecube}, \textsl{modfits}, \textsl{operate}, \textsl{readhead},
\textsl{sfinder}, \textsl{smeasure}, \textsl{splitcube}, and the library \textsl{tools}; and \emph{the author} refers to Alexander 
Men'shchikov of the D{\'e}partement d’Astrophysique, IRFU, CEA Saclay, France.

\vspace{0.5mm} 
It is assumed that anyone using this software (\emph{the user}) has read, understood, and agreed to the following conditions of use:

\vspace{0.5mm}
1. This software is freely available for download and use. A user is free to share this software with co-workers and students, but
when it is requested by a colleague not working directly with the user, the latter is asked to redirect that person to the
\textsl{getsf} web site: \url{http://irfu.cea.fr/Pisp/alexander.menshchikov/} or to the author:
\href{mailto:alexander.menshchikov@cea.fr}{\texttt{alexander.menshchikov@cea.fr}}

\vspace{0.5mm}
2. This software shall be used exclusively for research, education, non-profit, and non-military purposes. Specific written
permission from the author must be obtained before any commercial use of this software may be undertaken.

\vspace{0.5mm}
3. The banners of this software, these conditions of use, and information about the author shall remain with this software and
any descendent developed from and still based substantially upon this software.

\vspace{0.5mm}
4. The name of the institution with which the author is affiliated shall not be used to publicize any data and/or results
generated by this software. All findings and their interpretation are the opinions of the user and do not necessarily reflect those
of the author nor the institutions with which the author is affiliated.

\vspace{0.5mm} 
5. The author makes no representations about the suitability of this software for any purpose. Subject to the above conditions, 
this software is provided \emph{as is} without any expressed or implied warranty.

\section{Citations}
\label{citations}

The algorithms used in this software are described in:
\cite{Men'shchikov2021getsf,Men'shchikov2017,Men'shchikov2016,Men'shchikov2013,
Palmeirim_etal2013,Men'shchikov_etal2012,Men'shchikov_etal2010}. Detailed benchmarking of \textsl{getsf} and comparisons with its
predecessor \textsl{getsources} \citep{Men'shchikov_etal2012} is published by \cite{Men'shchikov2021bench}. Any publication
reporting results obtained using this software or any of its derivatives should include a citation of \cite{Men'shchikov2021getsf}
or a phrase like this: ``Use of \textsl{getsf}, developed by A. Men'shchikov at the DAp, IRFU, CEA Saclay, France, is hereby
acknowledged'' (if only individual scripts or utilities were used, then \textsl{getsf} may be substituted with their names).

%||||||||||||||||||||||||||||||||||||||||||||||||||||||||||||||||||||||||||||||||||||||||||||||||||||||||||||||||||||||||||||||||||
%||||||||||||||||||||||||||||||||||||||||||||||||||||||||||||||||||||||||||||||||||||||||||||||||||||||||||||||||||||||||||||||||||

\bibliographystyle{aa}
\bibliography{aamnem99,getsf.guide}

\end{document}